
\Data\ID{1003084901}
\Updated{2021-04-29 11:04:50}
\Level{A}
\SubArea{Properties of Sequences}
\LANGen{
\Question{
We were given a sequence \( \left( 2n \right)^{\infty}_{n=1} \). What does this formula express?}
{a sequence of all even natural numbers}{a sequence of all natural numbers}{a sequence of all odd natural numbers}{a sequence of all natural numbers that are divisible by \( 5 \)}{}{}}
\LANGpl{
\Question{
Dany jest ciąg \( \left( 2n \right)^{\infty}_{n=1} \). Co wyraża ten wzór?}
{ciąg wszystkich parzystych liczb naturalnych}{ciąg wszystkich liczb naturalnych}{ciąg wszystkich nieparzystych liczb naturalnych}{ciąg wszystkich liczb naturalnych podzielnych przez \( 5 \)}{}{}}
\LANGcs{
\Question{
Je dána posloupnost \( \left( 2n \right)^{\infty}_{n=1} \). Tento zápis vyjadřuje:}
{posloupnost všech sudých přirozených čísel}{posloupnost všech přirozených čísel}{posloupnost všech lichých přirozených čísel}{posloupnost všech přirozených čísel dělitelných pěti}{}{}}
\LANGsk{
\Question{
Je daná postupnosť \( \left( 2n \right)^{\infty}_{n=1} \). Tento zápis vyjadruje:}
{postupnosť všetkých párnych prirodzených čísel}{postupnosť všetkých prirodzených čísel}{postupnosť všetkých nepárnych prirodzených čísel}{postupnosť všetkých prirodzených čísel deliteľných piatimi}{}{}}
\LANGes{
\Question{
Dada la sucesión \( \left( 2n \right)^{\infty}_{n=1} \). ¿El conjunto de elementos de la sucesión es?}
{La secuencia de todos lo números naturales pares}{La secuencia de todos los números naturales}{La secuencia de todos lo números naturales impares}{La secuencia de todos los números naturales que son divisibles por \( 5 \)}{}{}}
\End

\Data\ID{1003084902}
\Updated{2022-04-23 05:04:02}
\Level{A}
\SubArea{Properties of Sequences}
\LANGen{
\Question{
We are given a sequence \( \left( 3n-2\right)^{\infty}_{n=1} \). 
What does this formula express?}
{a sequence of all natural numbers which after dividing by \( 3 \) give the remainder \( 1 \)}{a sequence of all natural numbers which are divisible by \( 3 \)}{a sequence of all natural numbers which are divisible by \( 2 \)}{a sequence of all odd natural numbers}{}{}}
\LANGpl{
\Question{
Dany jest ciąg  \( \left( 3n-2\right)^{\infty}_{n=1} \). 
Co wyraża ten wzór?}
{ciąg wszystkich liczb naturalnych, które po podzieleniu przez \( 3 \) dają resztę  \( 1 \)}{ciąg wszystkich liczb naturalnych podzielnych przez \( 3 \)}{ciąg wszystkich liczb naturalnych podzielnych przez \( 2 \)}{ciąg wszystkich nieparzystych liczb naturalnych}{}{}}
\LANGcs{
\Question{
Je dána posloupnost \( \left( 3n-2\right)^{\infty}_{n=1} \). 
Tento zápis vyjadřuje:}
{posloupnost všech přirozených čísel, která při dělení \( 3 \) dávají zbytek \( 1 \)}{posloupnost všech přirozených čísel dělitelných třemi}{posloupnost všech přirozených čísel dělitelných dvěma}{posloupnost všech lichých přirozených čísel}{}{}}
\LANGsk{
\Question{
Je daná postupnosť \( \left( 3n-2\right)^{\infty}_{n=1} \). 
Tento zápis vyjadruje:}
{postupnosť všetkých prirodzených čísel, ktoré pri delení \( 3 \) dávajú zvyšok \( 1 \)}{postupnosť všetkých prirodzených čísel deliteľných tromi}{postupnosť všetkých prirodzených čísel deliteľných dvoma}{postupnosť všetkých nepárnych prirodzených čísel}{}{}}
\LANGes{
\Question{
Dada la sucesión \( \left( 3n-2\right)^{\infty}_{n=1} \). 
¿Qué define esta fórmula?}
{La secuencia de todos los números naturales que después de dividirlos por \( 3 \) dan el resto \( 1 \)}{La secuencia de todos los números naturales que son divisibles por \( 3 \)}{La secuencia de todos los números naturales que son divisibles por \( 2 \)}{La secuencia de todos lo números naturales impares}{}{}}
\End

\Data\ID{1003084903}
\Updated{2021-05-03 07:05:00}
\Level{A}
\SubArea{Properties of Sequences}
\LANGen{
\Question{
The table contains ordered pairs of numbers \( [n;a_n] \). 
\[
\begin{array}{|c|c|c|c|c|c|} \hline n & 1 & 2 & 3 & 4 & 5 \\\hline a_n & -1 & 1 & -2 & 2 & -3 \\\hline
\end{array} 
\]
Which sequence is defined by the table?}
{\( \left(a_n\right)^5_{n=1}=-1\text{, }\ 1\text{, }-2\text{, }\ 2\text{, }-3 \)}{\( \left(a_n\right)^{10}_{n=1}=1\text{, }-1\text{, }\ 2\text{, }\ 1\text{, }\ 3\text{, }-2\text{, }\ 4\text{, }\ 2\text{, }\ 5\text{, }-3 \)}{\( \left(a_n\right)^5_{n=1}=1\text{, }\ 2\text{, }\ 3\text{, }\ 4\text{, }\ 5 \)}{\( \left(a_n\right)^5_{n=1}=0\text{, }\ 3\text{, }\ 1\text{, }\ 6\text{, }\ 2 \)}{}{}}
\LANGpl{
\Question{
Tabela zawiera uporządkowane ary liczb \( [n;a_n] \). 
\[
\begin{array}{|c|c|c|c|c|c|} \hline n & 1 & 2 & 3 & 4 & 5 \\\hline a_n & -1 & 1 & -2 & 2 & -3 \\\hline
\end{array} 
\]
Który ciąg określa ta tabela?}
{\( \left(a_n\right)^5_{n=1}=-1\text{, }\ 1\text{, }-2\text{, }\ 2\text{, }-3 \)}{\( \left(a_n\right)^{10}_{n=1}=1\text{, }-1\text{, }\ 2\text{, }\ 1\text{, }\ 3\text{, }-2\text{, }\ 4\text{, }\ 2\text{, }\ 5\text{, }-3 \)}{\( \left(a_n\right)^5_{n=1}=1\text{, }\ 2\text{, }\ 3\text{, }\ 4\text{, }\ 5 \)}{\( \left(a_n\right)^5_{n=1}=0\text{, }\ 3\text{, }\ 1\text{, }\ 6\text{, }\ 2 \)}{}{}}
\LANGcs{
\Question{
Uspořádané dvojice čísel \( [n;a_n] \) jsou zapsány do tabulky. 
\[
\begin{array}{|c|c|c|c|c|c|} \hline n & 1 & 2 & 3 & 4 & 5 \\\hline a_n & -1 & 1 & -2 & 2 & -3 \\\hline
\end{array} 
\]
Touto tabulkou je dána posloupnost:}
{\( \left(a_n\right)^5_{n=1}=-1\text{, }\ 1\text{, }-2\text{, }\ 2\text{, }-3 \)}{\( \left(a_n\right)^{10}_{n=1}=1\text{, }-1\text{, }\ 2\text{, }\ 1\text{, }\ 3\text{, }-2\text{, }\ 4\text{, }\ 2\text{, }\ 5\text{, }-3 \)}{\( \left(a_n\right)^5_{n=1}=1\text{, }\ 2\text{, }\ 3\text{, }\ 4\text{, }\ 5 \)}{\( \left(a_n\right)^5_{n=1}=0\text{, }\ 3\text{, }\ 1\text{, }\ 6\text{, }\ 2 \)}{}{}}
\LANGsk{
\Question{
Usporiadané dvojice čísel \( [n;a_n] \) sú zapísané do tabuľky. 
\[
\begin{array}{|c|c|c|c|c|c|} \hline n & 1 & 2 & 3 & 4 & 5 \\\hline a_n & -1 & 1 & -2 & 2 & -3 \\\hline
\end{array} 
\]
Touto tabuľkou je daná postupnosť:}
{\( \left(a_n\right)^5_{n=1}=-1\text{, }\ 1\text{, }-2\text{, }\ 2\text{, }-3 \)}{\( \left(a_n\right)^{10}_{n=1}=1\text{, }-1\text{, }\ 2\text{, }\ 1\text{, }\ 3\text{, }-2\text{, }\ 4\text{, }\ 2\text{, }\ 5\text{, }-3 \)}{\( \left(a_n\right)^5_{n=1}=1\text{, }\ 2\text{, }\ 3\text{, }\ 4\text{, }\ 5 \)}{\( \left(a_n\right)^5_{n=1}=0\text{, }\ 3\text{, }\ 1\text{, }\ 6\text{, }\ 2 \)}{}{}}
\LANGes{
\Question{
La tabla contiene pares ordenados de números \( [n;a_n] \). 
\[
\begin{array}{|c|c|c|c|c|c|} \hline n & 1 & 2 & 3 & 4 & 5 \\\hline a_n & -1 & 1 & -2 & 2 & -3 \\\hline
\end{array} 
\]
¿Qué sucesión define la tabla?}
{\( \left(a_n\right)^5_{n=1}=-1\text{, }\ 1\text{, }-2\text{, }\ 2\text{, }-3 \)}{\( \left(a_n\right)^{10}_{n=1}=1\text{, }-1\text{, }\ 2\text{, }\ 1\text{, }\ 3\text{, }-2\text{, }\ 4\text{, }\ 2\text{, }\ 5\text{, }-3 \)}{\( \left(a_n\right)^5_{n=1}=1\text{, }\ 2\text{, }\ 3\text{, }\ 4\text{, }\ 5 \)}{\( \left(a_n\right)^5_{n=1}=0\text{, }\ 3\text{, }\ 1\text{, }\ 6\text{, }\ 2 \)}{}{}}
\End

\Data\ID{1003084904}
\Updated{2021-05-03 07:05:32}
\Level{A}
\SubArea{Properties of Sequences}
\LANGen{
\Question{
We were given a sequence \( \left( \frac1n \right)^{\infty}_{n=1} \). 
Which expression describes the term \( a_{n-1}\ (n\in\mathbb{N}; n>1 ) \) of the given sequence?}
{\( a_{n-1} = \frac1{n-1} \)}{\( a_{n-1} = n-1 \)}{\( a_{n-1} = \frac1n \)}{\( a_{n-1} = \frac1{2n-1} \)}{}{}}
\LANGpl{
\Question{
Dany jest ciąg  \( \left( \frac1n \right)^{\infty}_{n=1} \). 
Które wyrażenie określa argument \( a_{n-1}\ (n\in\mathbb{N}; n>1 ) \) tego ciągu?}
{\( a_{n-1} = \frac1{n-1} \)}{\( a_{n-1} = n-1 \)}{\( a_{n-1} = \frac1n \)}{\( a_{n-1} = \frac1{2n-1} \)}{}{}}
\LANGcs{
\Question{
Je dána posloupnost \( \left( \frac1n \right)^{\infty}_{n=1} \). 
Pro člen \( a_{n-1}\ (n\in\mathbb{N}; n>1 ) \) této posloupnosti platí:}
{\( a_{n-1} = \frac1{n-1} \)}{\( a_{n-1} = n-1 \)}{\( a_{n-1} = \frac1n \)}{\( a_{n-1} = \frac1{2n-1} \)}{}{}}
\LANGsk{
\Question{
Je daná postupnosť \( \left( \frac1n \right)^{\infty}_{n=1} \). 
Pre člen \( a_{n-1}\ (n\in\mathbb{N}; n>1 ) \) tejto postupnosti platí:}
{\( a_{n-1} = \frac1{n-1} \)}{\( a_{n-1} = n-1 \)}{\( a_{n-1} = \frac1n \)}{\( a_{n-1} = \frac1{2n-1} \)}{}{}}
\LANGes{
\Question{
Dada la sucesión \( \left( \frac1n \right)^{\infty}_{n=1} \). 
Para el término \( a_{n-1}\ (n\in\mathbb{N}; n>1 ) \) de la sucesión es valida la igualdad:}
{\( a_{n-1} = \frac1{n-1} \)}{\( a_{n-1} = n-1 \)}{\( a_{n-1} = \frac1n \)}{\( a_{n-1} = \frac1{2n-1} \)}{}{}}
\End

\Data\ID{1003084906}
\Updated{2022-04-23 05:04:02}
\Level{A}
\SubArea{Properties of Sequences}
\LANGen{
\Question{
We are given a sequence \( \left( \frac{n+1}n \right)_{n=1}^{\infty} \). 
Which of the following formulations describes how is the given sequence defined?}
{defined by a formula for the \( n \)th term}{defined by a list of the sequence elements}{defined by a graph of the sequence}{defined by a recursive formula for the sequence}{}{}}
\LANGpl{
\Question{
Dany jest ciąg \( \left( \frac{n+1}n \right)_{n=1}^{\infty} \). 
Który z poniższych opisów definiuje ten ciąg?}
{wzór \( n \)-tego argumentu}{lista elementów ciągu}{wykres ciągu}{wzór rekurencyjny ciągu}{}{}}
\LANGcs{
\Question{
Je dána posloupnost \( \left( \frac{n+1}n \right)_{n=1}^{\infty} \). 
Vyberte možnost, která co nejlépe popisuje způsob zadání této posloupnosti.}
{vzorec pro \( n \)-tý člen}{výčet členů posloupnosti}{graf posloupnosti}{rekurentní vyjádření posloupnosti}{}{}}
\LANGsk{
\Question{
Je daná postupnosť \( \left( \frac{n+1}n \right)_{n=1}^{\infty} \). 
Vyberte možnosť, ktorá čo najlepšie popisuje spôsob zadania tejto postupnosti.}
{vzorec pre \( n \)-tý člen}{výber členov postupnosti}{graf postupnosti}{rekurentné vyjadrenie postupnosti}{}{}}
\LANGes{
\Question{
Dada la sucesión \( \left( \frac{n+1}n \right)_{n=1}^{\infty} \). 
¿Cuál de las siguientes opciones se ha usado para definirla?}
{la fórmula de  sus \( n \) términos}{una lista de términos de la sucesión}{un gráfico de la sucesión}{una fórmula recursiva de la sucesión}{}{}}
\End

\Data\ID{1003084907}
\Updated{2020-07-15 03:07:21}
\Level{A}
\SubArea{Properties of Sequences}
\LANGen{
\Question{
A sequence \( \left( a_n \right)^{\infty}_{n=1} \) is defined by the relations: \( a_1=3;\ a_{n+1}=\frac{a_n}{n+2}\text{, }n\in\mathbb{N} \). 
Which of the following descriptions is used to define the given sequence?}
{a recursion formula of a sequence}{a formula of the \(n\)th term}{a list of sequence elements}{a graph of a sequence}{}{}}
\LANGpl{
\Question{
Ciąg \( \left( a_n \right)^{\infty}_{n=1} \) określony jest następującymi zależnościami: \( a_1=3;\ a_{n+1}=\frac{a_n}{n+2}\text{, }n\in\mathbb{N} \). 
Który z poniższych opisów definiuje ten ciąg?}
{wzór rekurencyjny ciągu}{wzór  \(n\)-tego wyrazu tego ciągu}{lista wyrazów ciągu}{wykres ciągu}{}{}}
\LANGcs{
\Question{
Posloupnost \( \left( a_n \right)^{\infty}_{n=1} \) je dána vztahy: \( a_1=3;\ a_{n+1}=\frac{a_n}{n+2}\text{, }n\in\mathbb{N} \). 
Vyberte možnost, která co nejlépe popisuje způsob zadání této posloupnosti.}
{rekurentní vyjádření posloupnosti}{vzorec pro \(n\)-tý člen}{výčet členů posloupnosti}{graf posloupnosti}{}{}}
\LANGsk{
\Question{
Postupnosť \( \left( a_n \right)^{\infty}_{n=1} \) je daná vzťahmi: \( a_1=3;\ a_{n+1}=\frac{a_n}{n+2}\text{, }n\in\mathbb{N} \). 
Vyberte možnosť, ktorá čo najlepšie popisuje spôsob zadania tejto postupnosti.}
{rekurentné vyjadrenie postupnosti}{vzorec pre \(n\)-tý člen}{výber členov postupnosti}{graf postupnosti}{}{}}
\LANGes{
\Question{
La sucesión \( \left( a_n \right)^{\infty}_{n=1} \) está definida por las relaciones: \( a_1=3;\ a_{n+1}=\frac{a_n}{n+2}\text{, }n\in\mathbb{N} \). 
¿Cuál de las siguientes opciones se usa para definir la sucesión dada?}
{una fórmula recursiva de una sucesión}{la fórmula de \(n\) términos}{una lista de términos de sucesión}{un gráfico de sucesión}{}{}}
\End

\Data\ID{1003085001}
\Updated{2020-07-15 03:07:21}
\Level{A}
\SubArea{Properties of Sequences}
\LANGen{
\Question{
We were given the  sequence  \( \left(\frac1{3^n}\right)_{n=1}^{\infty} \).
What are its first five terms?}
{\( \frac13 \), \( \frac19 \), \( \frac1{27} \), \( \frac1{81} \), \( \frac1{243} \)}{\( 3 \), \( 9 \), \( 27 \), \( 81 \), \( 243 \)}{\( 3 \), \( 6 \), \( 9 \), \( 12 \), \( 15 \)}{\( \frac13 \), \( \frac16 \), \( \frac19 \), \( \frac1{12} \), \( \frac1{15} \)}{}{}}
\LANGpl{
\Question{
Dany jest ciąg  \( \left(\frac1{3^n}\right)_{n=1}^{\infty} \).
Podaj pięć pierwszych wyrazów tego ciągu?}
{\( \frac13 \), \( \frac19 \), \( \frac1{27} \), \( \frac1{81} \), \( \frac1{243} \)}{\( 3 \), \( 9 \), \( 27 \), \( 81 \), \( 243 \)}{\( 3 \), \( 6 \), \( 9 \), \( 12 \), \( 15 \)}{\( \frac13 \), \( \frac16 \), \( \frac19 \), \( \frac1{12} \), \( \frac1{15} \)}{}{}}
\LANGcs{
\Question{
Je dána posloupnost  \( \left(\frac1{3^n}\right)_{n=1}^{\infty} \).
Prvních pět členů této posloupnosti je:}
{\( \frac13 \), \( \frac19 \), \( \frac1{27} \), \( \frac1{81} \), \( \frac1{243} \)}{\( 3 \), \( 9 \), \( 27 \), \( 81 \), \( 243 \)}{\( 3 \), \( 6 \), \( 9 \), \( 12 \), \( 15 \)}{\( \frac13 \), \( \frac16 \), \( \frac19 \), \( \frac1{12} \), \( \frac1{15} \)}{}{}}
\LANGsk{
\Question{
Je daná postupnosť  \( \left(\frac1{3^n}\right)_{n=1}^{\infty} \).
Prvých päť členov tejto postupnosti je:}
{\( \frac13 \), \( \frac19 \), \( \frac1{27} \), \( \frac1{81} \), \( \frac1{243} \)}{\( 3 \), \( 9 \), \( 27 \), \( 81 \), \( 243 \)}{\( 3 \), \( 6 \), \( 9 \), \( 12 \), \( 15 \)}{\( \frac13 \), \( \frac16 \), \( \frac19 \), \( \frac1{12} \), \( \frac1{15} \)}{}{}}
\LANGes{
\Question{
Dada la sucesión \( \left(\frac1{3^n}\right)_{n=1}^{\infty} \).
Los primeros cinco términos son:}
{\( \frac13 \), \( \frac19 \), \( \frac1{27} \), \( \frac1{81} \), \( \frac1{243} \)}{\( 3 \), \( 9 \), \( 27 \), \( 81 \), \( 243 \)}{\( 3 \), \( 6 \), \( 9 \), \( 12 \), \( 15 \)}{\( \frac13 \), \( \frac16 \), \( \frac19 \), \( \frac1{12} \), \( \frac1{15} \)}{}{}}
\End

\Data\ID{1003085002}
\Updated{2020-07-15 03:07:21}
\Level{A}
\SubArea{Properties of Sequences}
\LANGen{
\Question{
We were given  the sequence  \( \left(\frac{n+3}{2n}\right)_{n=1}^{\infty} \).
What are its first five terms?}
{\( 2 \), \( \frac54 \), \( 1 \), \( \frac78 \), \( \frac45 \)}{\( \frac45 \), \( \frac78 \), \( 1 \), \( \frac54 \), \( 2 \)}{\( 2 \), \( \frac45 \), \( 1 \), \( \frac87 \), \( \frac54 \)}{\( \frac12 \), \( \frac23 \), \( \frac34 \), \( \frac45 \), \( \frac56 \)}{}{}}
\LANGpl{
\Question{
Dany jest ciąg  \( \left(\frac{n+3}{2n}\right)_{n=1}^{\infty} \).
Podaj pięć pierwszych wyrazów tego ciągu?}
{\( 2 \), \( \frac54 \), \( 1 \), \( \frac78 \), \( \frac45 \)}{\( \frac45 \), \( \frac78 \), \( 1 \), \( \frac54 \), \( 2 \)}{\( 2 \), \( \frac45 \), \( 1 \), \( \frac87 \), \( \frac54 \)}{\( \frac12 \), \( \frac23 \), \( \frac34 \), \( \frac45 \), \( \frac56 \)}{}{}}
\LANGcs{
\Question{
Je dána posloupnost  \( \left(\frac{n+3}{2n}\right)_{n=1}^{\infty} \).
Prvních pět členů této posloupnosti je:}
{\( 2 \), \( \frac54 \), \( 1 \), \( \frac78 \), \( \frac45 \)}{\( \frac45 \), \( \frac78 \), \( 1 \), \( \frac54 \), \( 2 \)}{\( 2 \), \( \frac45 \), \( 1 \), \( \frac87 \), \( \frac54 \)}{\( \frac12 \), \( \frac23 \), \( \frac34 \), \( \frac45 \), \( \frac56 \)}{}{}}
\LANGsk{
\Question{
Je daná postupnosť  \( \left(\frac{n+3}{2n}\right)_{n=1}^{\infty} \).
Prvých päť členov tejto postupnosti je:}
{\( 2 \), \( \frac54 \), \( 1 \), \( \frac78 \), \( \frac45 \)}{\( \frac45 \), \( \frac78 \), \( 1 \), \( \frac54 \), \( 2 \)}{\( 2 \), \( \frac45 \), \( 1 \), \( \frac87 \), \( \frac54 \)}{\( \frac12 \), \( \frac23 \), \( \frac34 \), \( \frac45 \), \( \frac56 \)}{}{}}
\LANGes{
\Question{
Dada la sucesión \( \left(\frac{n+3}{2n}\right)_{n=1}^{\infty} \).
Los primeros cinco términos son:}
{\( 2 \), \( \frac54 \), \( 1 \), \( \frac78 \), \( \frac45 \)}{\( \frac45 \), \( \frac78 \), \( 1 \), \( \frac54 \), \( 2 \)}{\( 2 \), \( \frac45 \), \( 1 \), \( \frac87 \), \( \frac54 \)}{\( \frac12 \), \( \frac23 \), \( \frac34 \), \( \frac45 \), \( \frac56 \)}{}{}}
\End

\Data\ID{1003085003}
\Updated{2020-07-15 03:07:21}
\Level{A}
\SubArea{Properties of Sequences}
\LANGen{
\Question{
We were given  the sequence \( \left(\sin\left(n\cdot\frac{\pi}2\right)\right)_{n=1}^{\infty} \).
What are the first five terms?}
{\( 1 \), \( 0 \), \( -1 \), \( 0 \), \( 1 \)}{\( 1 \), \( 0 \), \( 1 \), \( 0 \), \( 1 \)}{\( -1 \), \( 0 \), \( 1 \), \( 0 \), \( 1 \)}{\( 0 \), \( -1 \), \( 0 \), \( 1 \), \( 0 \)}{}{}}
\LANGpl{
\Question{
Dany jest ciąg  \( \left(\sin\left(n\cdot\frac{\pi}2\right)\right)_{n=1}^{\infty} \).
Podaj pięć pierwszych wyrazów tego ciągu?}
{\( 1 \), \( 0 \), \( -1 \), \( 0 \), \( 1 \)}{\( 1 \), \( 0 \), \( 1 \), \( 0 \), \( 1 \)}{\( -1 \), \( 0 \), \( 1 \), \( 0 \), \( 1 \)}{\( 0 \), \( -1 \), \( 0 \), \( 1 \), \( 0 \)}{}{}}
\LANGcs{
\Question{
Je dána posloupnost \( \left(\sin\left(n\cdot\frac{\pi}2\right)\right)_{n=1}^{\infty} \).
Prvních pět členů této posloupnosti je:}
{\( 1 \), \( 0 \), \( -1 \), \( 0 \), \( 1 \)}{\( 1 \), \( 0 \), \( 1 \), \( 0 \), \( 1 \)}{\( -1 \), \( 0 \), \( 1 \), \( 0 \), \( 1 \)}{\( 0 \), \( -1 \), \( 0 \), \( 1 \), \( 0 \)}{}{}}
\LANGsk{
\Question{
Je daná postupnosť \( \left(\sin\left(n\cdot\frac{\pi}2\right)\right)_{n=1}^{\infty} \).
Prvých päť členov tejto postupnosti je:}
{\( 1 \), \( 0 \), \( -1 \), \( 0 \), \( 1 \)}{\( 1 \), \( 0 \), \( 1 \), \( 0 \), \( 1 \)}{\( -1 \), \( 0 \), \( 1 \), \( 0 \), \( 1 \)}{\( 0 \), \( -1 \), \( 0 \), \( 1 \), \( 0 \)}{}{}}
\LANGes{
\Question{
Dada la sucesión \( \left(\sin\left(n\cdot\frac{\pi}2\right)\right)_{n=1}^{\infty} \).
Los primeros cinco términos son:}
{\( 1 \), \( 0 \), \( -1 \), \( 0 \), \( 1 \)}{\( 1 \), \( 0 \), \( 1 \), \( 0 \), \( 1 \)}{\( -1 \), \( 0 \), \( 1 \), \( 0 \), \( 1 \)}{\( 0 \), \( -1 \), \( 0 \), \( 1 \), \( 0 \)}{}{}}
\End

\Data\ID{1003085004}
\Updated{2020-07-15 03:07:21}
\Level{A}
\SubArea{Properties of Sequences}
\LANGen{
\Question{
A sequence \( \left(a_n\right)_{n=1}^{\infty} \) is defined by the  recursive  formula  \( a_1=1;\ a_{n+1} = 3a_n\text{, }n\in\mathbb{N} \).
What are its first five terms?}
{\( 1 \), \( 3 \), \( 9 \), \( 27 \), \( 81 \)}{\( 3 \), \( 9 \), \( 27 \), \( 81 \), \( 243 \)}{\( 1 \), \( 3 \), \( 6 \), \( 12 \), \( 24 \)}{\( 1 \), \( 3 \), \( 9 \), \( 30 \), \( 90 \)}{}{}}
\LANGpl{
\Question{
Ciąg  \( \left(a_n\right)_{n=1}^{\infty} \) określony jest wzorem rekurencyjnym  \( a_1=1;\ a_{n+1} = 3a_n\text{, }n\in\mathbb{N} \).
Podaj pięć pierwszych wyrazów tego ciągu?}
{\( 1 \), \( 3 \), \( 9 \), \( 27 \), \( 81 \)}{\( 3 \), \( 9 \), \( 27 \), \( 81 \), \( 243 \)}{\( 1 \), \( 3 \), \( 6 \), \( 12 \), \( 24 \)}{\( 1 \), \( 3 \), \( 9 \), \( 30 \), \( 90 \)}{}{}}
\LANGcs{
\Question{
Posloupnost  \( \left(a_n\right)_{n=1}^{\infty} \) je určena rekurentně:   \( a_1=1;\ a_{n+1} = 3a_n\text{, }n\in\mathbb{N} \).
Prvních pět členů této posloupnosti je:}
{\( 1 \), \( 3 \), \( 9 \), \( 27 \), \( 81 \)}{\( 3 \), \( 9 \), \( 27 \), \( 81 \), \( 243 \)}{\( 1 \), \( 3 \), \( 6 \), \( 12 \), \( 24 \)}{\( 1 \), \( 3 \), \( 9 \), \( 30 \), \( 90 \)}{}{}}
\LANGsk{
\Question{
Postupnosť  \( \left(a_n\right)_{n=1}^{\infty} \) je určená rekurentne:   \( a_1=1;\ a_{n+1} = 3a_n\text{, }n\in\mathbb{N} \).
Prvých päť členov tejto postupnosti je:}
{\( 1 \), \( 3 \), \( 9 \), \( 27 \), \( 81 \)}{\( 3 \), \( 9 \), \( 27 \), \( 81 \), \( 243 \)}{\( 1 \), \( 3 \), \( 6 \), \( 12 \), \( 24 \)}{\( 1 \), \( 3 \), \( 9 \), \( 30 \), \( 90 \)}{}{}}
\LANGes{
\Question{
La sucesión \( \left(a_n\right)_{n=1}^{\infty} \) viene dada por la fórmula recursiva  \( a_1=1;\ a_{n+1} = 3a_n\text{, }n\in\mathbb{N} \).
Los primeros cinco términos son:}
{\( 1 \), \( 3 \), \( 9 \), \( 27 \), \( 81 \)}{\( 3 \), \( 9 \), \( 27 \), \( 81 \), \( 243 \)}{\( 1 \), \( 3 \), \( 6 \), \( 12 \), \( 24 \)}{\( 1 \), \( 3 \), \( 9 \), \( 30 \), \( 90 \)}{}{}}
\End

\Data\ID{1003085005}
\Updated{2020-07-15 03:07:21}
\Level{A}
\SubArea{Properties of Sequences}
\LANGen{
\Question{
A sequence \( \left( a_n \right)_{n=1}^{\infty} \) is defined by  the  recursive  formula  \( a_1=1;\ a_{n+1}=\frac1{1+a_n}\text{, }n\in\mathbb{N} \).
What are its first five terms?}
{\( 1 \), \( \frac12 \), \( \frac23 \), \( \frac35 \), \( \frac58 \)}{\( 1 \), \( \frac12 \), \( \frac23 \), \( \frac34 \), \( \frac58 \)}{\( 1 \), \( 2 \), \( \frac32 \), \( \frac53 \), \( \frac85 \)}{\( 1 \), \( \frac12 \), \( \frac32 \), \( \frac35 \), \( \frac85 \)}{}{}}
\LANGpl{
\Question{
Ciąg \( \left( a_n \right)_{n=1}^{\infty} \) określony jest wzorem rekurencyjnym \( a_1=1;\ a_{n+1}=\frac1{1+a_n}\text{, }n\in\mathbb{N} \).
Podaj pięć pierwszych wyrazów tego ciągu?}
{\( 1 \), \( \frac12 \), \( \frac23 \), \( \frac35 \), \( \frac58 \)}{\( 1 \), \( \frac12 \), \( \frac23 \), \( \frac34 \), \( \frac58 \)}{\( 1 \), \( 2 \), \( \frac32 \), \( \frac53 \), \( \frac85 \)}{\( 1 \), \( \frac12 \), \( \frac32 \), \( \frac35 \), \( \frac85 \)}{}{}}
\LANGcs{
\Question{
Posloupnost \( \left( a_n \right)_{n=1}^{\infty} \) je určena rekurentně:  \( a_1=1;\ a_{n+1}=\frac1{1+a_n}\text{, }n\in\mathbb{N} \).
Prvních pět členů této posloupnosti je:}
{\( 1 \), \( \frac12 \), \( \frac23 \), \( \frac35 \), \( \frac58 \)}{\( 1 \), \( \frac12 \), \( \frac23 \), \( \frac34 \), \( \frac58 \)}{\( 1 \), \( 2 \), \( \frac32 \), \( \frac53 \), \( \frac85 \)}{\( 1 \), \( \frac12 \), \( \frac32 \), \( \frac35 \), \( \frac85 \)}{}{}}
\LANGsk{
\Question{
Postupnosť \( \left( a_n \right)_{n=1}^{\infty} \) je určená rekurentne:  \( a_1=1;\ a_{n+1}=\frac1{1+a_n}\text{, }n\in\mathbb{N} \).
Prvých päť členov tejto postupnosti je:}
{\( 1 \), \( \frac12 \), \( \frac23 \), \( \frac35 \), \( \frac58 \)}{\( 1 \), \( \frac12 \), \( \frac23 \), \( \frac34 \), \( \frac58 \)}{\( 1 \), \( 2 \), \( \frac32 \), \( \frac53 \), \( \frac85 \)}{\( 1 \), \( \frac12 \), \( \frac32 \), \( \frac35 \), \( \frac85 \)}{}{}}
\LANGes{
\Question{
La sucesión \( \left( a_n \right)_{n=1}^{\infty} \) viene dada por la fórmula recursiva  \( a_1=1;\ a_{n+1}=\frac1{1+a_n}\text{, }n\in\mathbb{N} \).
Los primeros cinco términos son:}
{\( 1 \), \( \frac12 \), \( \frac23 \), \( \frac35 \), \( \frac58 \)}{\( 1 \), \( \frac12 \), \( \frac23 \), \( \frac34 \), \( \frac58 \)}{\( 1 \), \( 2 \), \( \frac32 \), \( \frac53 \), \( \frac85 \)}{\( 1 \), \( \frac12 \), \( \frac32 \), \( \frac35 \), \( \frac85 \)}{}{}}
\End

\Data\ID{1003085006}
\Updated{2020-07-15 03:07:21}
\Level{A}
\SubArea{Properties of Sequences}
\LANGen{
\Question{
A sequence \( \left(a_n\right)_{n=1}^{\infty} \)  is defined by  the recursive  formula \( a_1=1\text{, }a_2=2;\ a_{n+2} = \frac12\left(a_{n+1}+a_n\right)\text{, }n\in\mathbb{N} \).
What are its first five terms?}
{\( 1 \), \( 2 \), \( \frac32 \), \( \frac74 \), \( \frac{13}8 \)}{\( 1 \), \( 2 \), \( \frac32 \), \( \frac47 \), \( \frac8{13} \)}{\( 1 \), \( 2 \), \( 3 \), \( 7 \), \( 13 \)}{\( 1 \), \( 2 \), \( \frac23 \), \( \frac47 \), \( \frac{13}8 \)}{}{}}
\LANGpl{
\Question{
Ciąg \( \left(a_n\right)_{n=1}^{\infty} \)  określony jest wzorem rekurencyjnym \( a_1=1\text{, }a_2=2;\ a_{n+2} = \frac12\left(a_{n+1}+a_n\right)\text{, }n\in\mathbb{N} \).
Podaj pięć pierwszych wyrazów tego ciągu?}
{\( 1 \), \( 2 \), \( \frac32 \), \( \frac74 \), \( \frac{13}8 \)}{\( 1 \), \( 2 \), \( \frac32 \), \( \frac47 \), \( \frac8{13} \)}{\( 1 \), \( 2 \), \( 3 \), \( 7 \), \( 13 \)}{\( 1 \), \( 2 \), \( \frac23 \), \( \frac47 \), \( \frac{13}8 \)}{}{}}
\LANGcs{
\Question{
Posloupnost \( \left(a_n\right)_{n=1}^{\infty} \)  je určena rekurentně: \( a_1=1\text{, }a_2=2;\ a_{n+2} = \frac12\left(a_{n+1}+a_n\right)\text{, }n\in\mathbb{N} \).
Prvních pět členů této posloupnosti je:}
{\( 1 \), \( 2 \), \( \frac32 \), \( \frac74 \), \( \frac{13}8 \)}{\( 1 \), \( 2 \), \( \frac32 \), \( \frac47 \), \( \frac8{13} \)}{\( 1 \), \( 2 \), \( 3 \), \( 7 \), \( 13 \)}{\( 1 \), \( 2 \), \( \frac23 \), \( \frac47 \), \( \frac{13}8 \)}{}{}}
\LANGsk{
\Question{
Postupnosť \( \left(a_n\right)_{n=1}^{\infty} \)  je určená rekurentne: \( a_1=1\text{, }a_2=2;\ a_{n+2} = \frac12\left(a_{n+1}+a_n\right)\text{, }n\in\mathbb{N} \).
Prvých päť členov tejto postupnosti je:}
{\( 1 \), \( 2 \), \( \frac32 \), \( \frac74 \), \( \frac{13}8 \)}{\( 1 \), \( 2 \), \( \frac32 \), \( \frac47 \), \( \frac8{13} \)}{\( 1 \), \( 2 \), \( 3 \), \( 7 \), \( 13 \)}{\( 1 \), \( 2 \), \( \frac23 \), \( \frac47 \), \( \frac{13}8 \)}{}{}}
\LANGes{
\Question{
La sucesión \( \left(a_n\right)_{n=1}^{\infty} \)  viene dada por la fórmula recursiva \( a_1=1\text{, }a_2=2;\ a_{n+2} = \frac12\left(a_{n+1}+a_n\right)\text{, }n\in\mathbb{N} \).
Los primeros cinco términos son:}
{\( 1 \), \( 2 \), \( \frac32 \), \( \frac74 \), \( \frac{13}8 \)}{\( 1 \), \( 2 \), \( \frac32 \), \( \frac47 \), \( \frac8{13} \)}{\( 1 \), \( 2 \), \( 3 \), \( 7 \), \( 13 \)}{\( 1 \), \( 2 \), \( \frac23 \), \( \frac47 \), \( \frac{13}8 \)}{}{}}
\End

\Data\ID{1003085007}
\Updated{2020-07-15 03:07:21}
\Level{A}
\SubArea{Properties of Sequences}
\LANGen{
\Question{
We were given the  sequence \( \left( (-1)^{n+1} \right)_{n=1}^{\infty} \).
What is its seventh term?}
{\( 1 \)}{\( -1 \)}{\( 0 \)}{\( 2 \)}{}{}}
\LANGpl{
\Question{
Dany jest ciąg \( \left( (-1)^{n+1} \right)_{n=1}^{\infty} \).
Jaki jest siódmy wyraz tego ciągu?}
{\( 1 \)}{\( -1 \)}{\( 0 \)}{\( 2 \)}{}{}}
\LANGcs{
\Question{
Je dána posloupnost \( \left( (-1)^{n+1} \right)_{n=1}^{\infty} \).
Sedmý člen této posloupnosti je:}
{\( 1 \)}{\( -1 \)}{\( 0 \)}{\( 2 \)}{}{}}
\LANGsk{
\Question{
Je daná postupnosť \( \left( (-1)^{n+1} \right)_{n=1}^{\infty} \).
Siedmy člen tejto postupnosti je:}
{\( 1 \)}{\( -1 \)}{\( 0 \)}{\( 2 \)}{}{}}
\LANGes{
\Question{
Dada la sucesión \( \left( (-1)^{n+1} \right)_{n=1}^{\infty} \).
El séptimo término de la sucesión es:}
{\( 1 \)}{\( -1 \)}{\( 0 \)}{\( 2 \)}{}{}}
\End

\Data\ID{1003085008}
\Updated{2020-07-15 03:07:21}
\Level{A}
\SubArea{Properties of Sequences}
\LANGen{
\Question{
A sequence \( \left( a_n \right)_{n=1}^{\infty} \) is defined by the  recursive  formula: \( a_1=1;\ a_{n+1}=-2a_n\text{, }n\in\mathbb{N} \).
What is its third term?}
{\( 4 \)}{\( 2 \)}{\( -4 \)}{\( \frac12 \)}{}{}}
\LANGpl{
\Question{
Ciąg \( \left( a_n \right)_{n=1}^{\infty} \) określony jest wzorem  rekurencyjnym: \( a_1=1;\ a_{n+1}=-2a_n\text{, }n\in\mathbb{N} \).
Jaki jest jego trzeci wyraz?}
{\( 4 \)}{\( 2 \)}{\( -4 \)}{\( \frac12 \)}{}{}}
\LANGcs{
\Question{
Posloupnost \( \left( a_n \right)_{n=1}^{\infty} \) je určena rekurentně: \( a_1=1;\ a_{n+1}=-2a_n\text{, }n\in\mathbb{N} \).
Třetí člen této posloupnosti je:}
{\( 4 \)}{\( 2 \)}{\( -4 \)}{\( \frac12 \)}{}{}}
\LANGsk{
\Question{
Postupnosť \( \left( a_n \right)_{n=1}^{\infty} \) je určená rekurentne: \( a_1=1;\ a_{n+1}=-2a_n\text{, }n\in\mathbb{N} \).
Tretí člen tejto postupnosti je:}
{\( 4 \)}{\( 2 \)}{\( -4 \)}{\( \frac12 \)}{}{}}
\LANGes{
\Question{
La sucesión \( \left( a_n \right)_{n=1}^{\infty} \) viene dada por la fórmula recursiva: \( a_1=1;\ a_{n+1}=-2a_n\text{, }n\in\mathbb{N} \).
El tercer término de la sucesión es:}
{\( 4 \)}{\( 2 \)}{\( -4 \)}{\( \frac12 \)}{}{}}
\End

\Data\ID{1003085009}
\Updated{2020-07-15 03:07:21}
\Level{A}
\SubArea{Properties of Sequences}
\LANGen{
\Question{
We were given the  sequence \( \left(\frac{n}{n+2}\right)_{n=1}^{\infty} \).
What is its tenth term?}
{\( \frac56 \)}{\( \frac65 \)}{\( 1 \)}{\( 10 \)}{}{}}
\LANGpl{
\Question{
Dany jest ciąg \( \left(\frac{n}{n+2}\right)_{n=1}^{\infty} \).
Jaki jest jego dziesiąty wyraz?}
{\( \frac56 \)}{\( \frac65 \)}{\( 1 \)}{\( 10 \)}{}{}}
\LANGcs{
\Question{
Je dána posloupnost \( \left(\frac{n}{n+2}\right)_{n=1}^{\infty} \).
Desátý člen této posloupnosti je:}
{\( \frac56 \)}{\( \frac65 \)}{\( 1 \)}{\( 10 \)}{}{}}
\LANGsk{
\Question{
Je daná postupnosť \( \left(\frac{n}{n+2}\right)_{n=1}^{\infty} \).
Desiaty člen tejto postupnosti je:}
{\( \frac56 \)}{\( \frac65 \)}{\( 1 \)}{\( 10 \)}{}{}}
\LANGes{
\Question{
Dada la sucesión \( \left(\frac{n}{n+2}\right)_{n=1}^{\infty} \).
El décimo término de la sucesión es:}
{\( \frac56 \)}{\( \frac65 \)}{\( 1 \)}{\( 10 \)}{}{}}
\End

\Data\ID{1003085010}
\Updated{2021-05-03 07:05:27}
\Level{A}
\SubArea{Properties of Sequences}
\LANGen{
\Question{
We were given two sequences \( \left( 2^{2n-2} \right)_{n=1}^{\infty} \) and \( \left( n^2 \right)_{n=1}^{\infty} \).
What  is the ratio of their fourth terms?}
{\( 4:1 \)}{\( 2:1 \)}{\( 3:1 \)}{\( 1:1 \)}{}{}}
\LANGpl{
\Question{
Dane są dwa ciągi  \( \left( 2^{2n-2} \right)_{n=1}^{\infty} \) i \( \left( n^2 \right)_{n=1}^{\infty} \).
Jaki jest stosunek ich czwartych wyrazów?}
{\( 4:1 \)}{\( 2:1 \)}{\( 3:1 \)}{\( 1:1 \)}{}{}}
\LANGcs{
\Question{
Jsou dány posloupnosti \( \left( 2^{2n-2} \right)_{n=1}^{\infty} \) a \( \left( n^2 \right)_{n=1}^{\infty} \).
V jakém poměru jsou čtvrté členy obou posloupností?}
{\( 4:1 \)}{\( 2:1 \)}{\( 3:1 \)}{\( 1:1 \)}{}{}}
\LANGsk{
\Question{
Sú dané postupnosti \( \left( 2^{2n-2} \right)_{n=1}^{\infty} \) a \( \left( n^2 \right)_{n=1}^{\infty} \).
V akom pomere sú štvrté členy oboch postupností?}
{\( 4:1 \)}{\( 2:1 \)}{\( 3:1 \)}{\( 1:1 \)}{}{}}
\LANGes{
\Question{
Dadas dos sucesiones \( \left( 2^{2n-2} \right)_{n=1}^{\infty} \) y \( \left( n^2 \right)_{n=1}^{\infty} \).
¿En qué proporción están los términos cuartos de estas sucesiones?}
{\( 4:1 \)}{\( 2:1 \)}{\( 3:1 \)}{\( 1:1 \)}{}{}}
\End

\Data\ID{1003107310}
\Updated{2020-07-15 03:07:12}
\Level{A}
\SubArea{Properties of Sequences}
\LANGen{
\Question{
We are given a sequence \( \left( a_n \right)^{\infty}_{n=1} \)   defined recursively by: \( a_1=1,\ a_2=2\,;\ a_{n+2}=\frac12\left( a_{n+1}+a_n\right),\ n\in\mathbb{N} \).
Find the sum of the first four terms of this sequence.}
{\( \frac{25}4 \)}{\( \frac{63}8 \)}{\( \frac{13}4 \)}{\( \frac4{25} \)}{}{}}
\LANGpl{
\Question{
Dany jest ciąg \( \left( a_n \right)^{\infty}_{n=1} \)   określony rekurencyjnie przez: \( a_1=1,\ a_2=2\,;\ a_{n+2}=\frac12\left( a_{n+1}+a_n\right),\ n\in\mathbb{N} \).
Wyznacz sumę czterech pierwszych wyrazów tego ciągu.}
{\( \frac{25}4 \)}{\( \frac{63}8 \)}{\( \frac{13}4 \)}{\( \frac4{25} \)}{}{}}
\LANGcs{
\Question{
Posloupnost \( \left( a_n \right)^{\infty}_{n=1} \)  je určena rekurentně: \( a_1=1,\ a_2=2\,;\ a_{n+2}=\frac12\left( a_{n+1}+a_n\right),\ n\in\mathbb{N} \).
Součet prvních čtyř členů této posloupnosti je roven:}
{\( \frac{25}4 \)}{\( \frac{63}8 \)}{\( \frac{13}4 \)}{\( \frac4{25} \)}{}{}}
\LANGsk{
\Question{
Postupnosť \( \left( a_n \right)^{\infty}_{n=1} \)  je určená rekurentne: \( a_1=1,\ a_2=2\,;\ a_{n+2}=\frac12\left( a_{n+1}+a_n\right),\ n\in\mathbb{N} \).
Súčet prvých štyroch členov tejto postupnosti je rovný:}
{\( \frac{25}4 \)}{\( \frac{63}8 \)}{\( \frac{13}4 \)}{\( \frac4{25} \)}{}{}}
\LANGes{
\Question{
La sucesión \( \left( a_n \right)^{\infty}_{n=1} \)  viene dada por la fórmula recursiva: \( a_1=1,\ a_2=2\,;\ a_{n+2}=\frac12\left( a_{n+1}+a_n\right),\ n\in\mathbb{N} \).
Halla la suma de los cuatro primeros términos de la sucesión.}
{\( \frac{25}4 \)}{\( \frac{63}8 \)}{\( \frac{13}4 \)}{\( \frac4{25} \)}{}{}}
\End

\Data\ID{1103084905}
\Updated{2022-04-23 05:04:02}
\Level{A}
\SubArea{Properties of Sequences}
\IMG{1103084905.pdf}
\LANGen{
\Question{
Which sequence is defined by the given graph?}
{\( \left( a_n \right)^{5}_{n=1} = 2\text{, }\ 1\text{, }\ 3\text{, }\ 1\text{, }\ 2 \)}{\( \left( a_n \right)^{10}_{n=1} = 1\text{, }\ 2\text{, }\ 2\text{, }\ 1\text{, }\ 3\text{, }\ 3\text{, }\ 4\text{, }\ 1\text{, }\ 5\text{, }\ 2 \)}{\( \left( a_n \right)^{5}_{n=1} = 1\text{, }\ 2\text{, }\ 3\text{, }\ 4\text{, }\ 5 \)}{\( \left( a_n \right)^{5}_{n=1} = 1\text{, }\ 1\text{, }\ 2\text{, }\ 2\text{, }\ 3 \)}{}{}}
\LANGpl{
\Question{
Który ciąg określony został przez poniższy wykres?}
{\( \left( a_n \right)^{5}_{n=1} = 2\text{, }\ 1\text{, }\ 3\text{, }\ 1\text{, }\ 2 \)}{\( \left( a_n \right)^{10}_{n=1} = 1\text{, }\ 2\text{, }\ 2\text{, }\ 1\text{, }\ 3\text{, }\ 3\text{, }\ 4\text{, }\ 1\text{, }\ 5\text{, }\ 2 \)}{\( \left( a_n \right)^{5}_{n=1} = 1\text{, }\ 2\text{, }\ 3\text{, }\ 4\text{, }\ 5 \)}{\( \left( a_n \right)^{5}_{n=1} = 1\text{, }\ 1\text{, }\ 2\text{, }\ 2\text{, }\ 3 \)}{}{}}
\LANGcs{
\Question{
Vyberte posloupnost, která je dána grafem na obrázku.}
{\( \left( a_n \right)^{5}_{n=1} = 2\text{, }\ 1\text{, }\ 3\text{, }\ 1\text{, }\ 2 \)}{\( \left( a_n \right)^{10}_{n=1} = 1\text{, }\ 2\text{, }\ 2\text{, }\ 1\text{, }\ 3\text{, }\ 3\text{, }\ 4\text{, }\ 1\text{, }\ 5\text{, }\ 2 \)}{\( \left( a_n \right)^{5}_{n=1} = 1\text{, }\ 2\text{, }\ 3\text{, }\ 4\text{, }\ 5 \)}{\( \left( a_n \right)^{5}_{n=1} = 1\text{, }\ 1\text{, }\ 2\text{, }\ 2\text{, }\ 3 \)}{}{}}
\LANGsk{
\Question{
Vyberte postupnosť, ktorá je daná grafom.}
{\( \left( a_n \right)^{5}_{n=1} = 2\text{, }\ 1\text{, }\ 3\text{, }\ 1\text{, }\ 2 \)}{\( \left( a_n \right)^{10}_{n=1} = 1\text{, }\ 2\text{, }\ 2\text{, }\ 1\text{, }\ 3\text{, }\ 3\text{, }\ 4\text{, }\ 1\text{, }\ 5\text{, }\ 2 \)}{\( \left( a_n \right)^{5}_{n=1} = 1\text{, }\ 2\text{, }\ 3\text{, }\ 4\text{, }\ 5 \)}{\( \left( a_n \right)^{5}_{n=1} = 1\text{, }\ 1\text{, }\ 2\text{, }\ 2\text{, }\ 3 \)}{}{}}
\LANGes{
\Question{
Elige la sucesión representada por el gráfico.}
{\( \left( a_n \right)^{5}_{n=1} = 2\text{, }\ 1\text{, }\ 3\text{, }\ 1\text{, }\ 2 \)}{\( \left( a_n \right)^{10}_{n=1} = 1\text{, }\ 2\text{, }\ 2\text{, }\ 1\text{, }\ 3\text{, }\ 3\text{, }\ 4\text{, }\ 1\text{, }\ 5\text{, }\ 2 \)}{\( \left( a_n \right)^{5}_{n=1} = 1\text{, }\ 2\text{, }\ 3\text{, }\ 4\text{, }\ 5 \)}{\( \left( a_n \right)^{5}_{n=1} = 1\text{, }\ 1\text{, }\ 2\text{, }\ 2\text{, }\ 3 \)}{}{}}
\End

\Data\ID{1103084908}
\Updated{2020-07-15 03:07:10}
\Level{A}
\SubArea{Properties of Sequences}
\LANGen{
\Question{
We were given a sequence \( \left( 1.5n\right)^{\infty}_{n=1} \). 
Which of the diagrams shows correctly the first five elements of the sequence?}
{}{}{}{}{}{}}
\LANGpl{
\Question{
Dany jest ciąg \( \left( 1{,}5n\right)^{\infty}_{n=1} \). 
Który z  diagramów poprawnie wskazuje pięć pierwszych wyrazów tego ciągu?}
{}{}{}{}{}{}}
\LANGcs{
\Question{
Je dána posloupnost \( \left( 1{,}5n\right)^{\infty}_{n=1} \). 
Vyberte obrázek, na kterém je správně vyznačeno prvních pět členů dané posloupnosti.}
{}{}{}{}{}{}}
\LANGsk{
\Question{
Je daná postupnosť \( \left( 1{,}5n\right)^{\infty}_{n=1} \). 
Vyberte obrázok, na ktorom je správne vyznačených prvých päť členov danej postupnosti.}
{}{}{}{}{}{}}
\LANGes{
\Question{
Dada la sucesión \( \left( 1.5n\right)^{\infty}_{n=1} \). 
Elige el diagrama que representa los cinco primeros términos de la sucesión.}
{}{}{}{}{}{}}

\IMGanswer{1103084908a.pdf}
\IMGanswer{1103084908b.pdf}
\IMGanswer{1103084908c.pdf}
\IMGanswer{1103084908d.pdf}\End

\Data\ID{2000005805}
\Updated{2022-08-25 10:08:29}
\Level{A}
\SubArea{Properties of Sequences}
\LANGen{
\Question{
In the sequence \( \{16,20,24,28,\dots\}\) each term can be obtained from previous term by adding \(4\). Which of the following numbers  is not a term in the sequence?}
{\(322\)}{\(200\)}{\(40\)}{\(764\)}{}{}}
\LANGpl{
\Question{
W ciągu \( \{16,20,24,28,\dots\}\) każdy wyraz można otrzymać z poprzedniego wyrazu przez dodanie \(4\). Która z poniższych liczb nie jest wyrazem w ciągu?}
{\(322\)}{\(200\)}{\(40\)}{\(764\)}{}{}}
\LANGcs{
\Question{
V posloupnosti \( \{16,20,24,28,\dots\}\) je každý další člen o \(4\) větší než předcházející. Které z daných čísel není členem této posloupnosti?}
{\(322\)}{\(200\)}{\(40\)}{\(764\)}{}{}}
\LANGsk{
\Question{
V postupnosti \( \{16,20,24,28,\dots\}\) je každý ďalší člen väčší o 4 ako predchádzajúci. Ktoré z daných čísel nemôže byť členom tejto postupnosti?}
{\(322\)}{\(200\)}{\(40\)}{\(764\)}{}{}}
\LANGes{
\Question{
En la sucesión \( \{16,20,24,28,\dots\}\) cada término se puede obtener del término anterior sumando un \(4\). ¿Cuál de los siguientes números no es uno de los términos de la sucesión dada?}
{\(322\)}{\(200\)}{\(40\)}{\(764\)}{}{}}
\End

\Data\ID{2000005806}
\Updated{2022-08-25 10:08:29}
\Level{A}
\SubArea{Properties of Sequences}
\LANGen{
\Question{
In the sequence \( \{-45,-36,-27,-18,\dots\}\) each term can be obtained from previous term by adding \(9\). Which of the following numbers  is not a term in the sequence?}
{\(285\)}{\(135\)}{\(927\)}{\(252\)}{}{}}
\LANGpl{
\Question{
W ciągu \( \{-45,-36,-27,-18,\dots\}\) każdy wyraz można otrzymać z poprzedniego wyrazu przez dodanie \(9\). Która z poniższych liczb nie jest wyrazem w ciągu?}
{\(285\)}{\(135\)}{\(927\)}{\(252\)}{}{}}
\LANGcs{
\Question{
V posloupnosti \( \{-45,-36,-27,-18,\dots\}\) je každý další člen o \(9\) větší než předcházející. Které z daných čísel není členem této poslouposti?}
{\(285\)}{\(135\)}{\(927\)}{\(252\)}{}{}}
\LANGsk{
\Question{
V postupnosti \( \{-45,-36,-27,-18,\dots\}\) je každý ďalší člen väčší o 9 ako predchádzajúci. Ktoré z daných čísel nemôže byť členom tejto postupnosti?}
{\(285\)}{\(135\)}{\(927\)}{\(252\)}{}{}}
\LANGes{
\Question{
En la sucesión \( \{-45,-36,-27,-18,\dots\}\) cada término se puede obtener del término anterior sumando un \(9\). ¿Cuál de los siguientes números no es uno de los términos de la sucesión dada?}
{\(285\)}{\(135\)}{\(927\)}{\(252\)}{}{}}
\End

\Data\ID{2000005807}
\Updated{2022-08-25 10:08:29}
\Level{A}
\SubArea{Properties of Sequences}
\LANGen{
\Question{
We are given a sequence with the same difference between consecutive terms. We know it starts at \(-7\) and its \(4\)th term is \(11\). Find its \(12\)th term.}
{\(59\)}{\(53\)}{\(65\)}{\(71\)}{}{}}
\LANGpl{
\Question{
Otrzymujemy ciąg z taką samą różnicą między kolejnymi wyrazami. Wiemy, że zaczyna się od \(-7\), a jego \(4\)-ty wyraz to \(11\). Znajdź jego \(12\) wyraz.}
{\(59\)}{\(53\)}{\(65\)}{\(71\)}{}{}}
\LANGcs{
\Question{
Je dána posloupnost se stejným rozdílem dvou následujících členů. Posloupnost začíná v \(-7\) a \(4\). člen je \(11\). Určete \(12\). člen.}
{\(59\)}{\(53\)}{\(65\)}{\(71\)}{}{}}
\LANGsk{
\Question{
Daná je postupnosť s rovnakým rozdielom medzi nasledujúcimi členmi. Postupnosť začína v \(-7\) a  \(4\). člen je \(11\). Určte dvanásty člen.}
{\(59\)}{\(53\)}{\(65\)}{\(71\)}{}{}}
\LANGes{
\Question{
Dada una sucesión con la misma diferencia entre términos consecutivos. Sabemos que comienza en \(-7\) y su cuarto término es \(11\). Calcula su duodécimo término.}
{\(59\)}{\(53\)}{\(65\)}{\(71\)}{}{}}
\End

\Data\ID{2000005809}
\Updated{2022-08-25 10:08:29}
\Level{A}
\SubArea{Properties of Sequences}
\LANGen{
\Question{
Which of the following sequences does not contain a term equal to \(256\)?}
{\(a_n=(-1)^n(n+1)^2\)}{\(a_n=\frac{(-2)^n}{n+4}\)}{\(a_n = \frac{n^2-516}{n}\)}{\( a_n=(-1)^n2^n\)}{}{}}
\LANGpl{
\Question{
Który z poniższych ciągów nie zawiera wyrażenia równego \(256\)?}
{\(a_n=(-1)^n(n+1)^2\)}{\(a_n=\frac{(-2)^n}{n+4}\)}{\(a_n = \frac{n^2-516}{n}\)}{\( a_n=(-1)^n2^n\)}{}{}}
\LANGcs{
\Question{
Která z daných posloupností neobsahuje člen \(256\)?}
{\(a_n=(-1)^n(n+1)^2\)}{\(a_n=\frac{(-2)^n}{n+4}\)}{\(a_n = \frac{n^2-516}{n}\)}{\( a_n=(-1)^n2^n\)}{}{}}
\LANGsk{
\Question{
Ktorá z nasledujúcich postupností neobsahuje člen \(256\)?}
{\(a_n=(-1)^n(n+1)^2\)}{\(a_n=\frac{(-2)^n}{n+4}\)}{\(a_n = \frac{n^2-516}{n}\)}{\( a_n=(-1)^n2^n\)}{}{}}
\LANGes{
\Question{
¿Cuál de las siguientes sucesiones no contiene un término equivalente a \(256\)?}
{\(a_n=(-1)^n(n+1)^2\)}{\(a_n=\frac{(-2)^n}{n+4}\)}{\(a_n = \frac{n^2-516}{n}\)}{\( a_n=(-1)^n2^n\)}{}{}}
\End

\Data\ID{2000005810}
\Updated{2023-03-27 11:03:25}
\Level{A}
\SubArea{Properties of Sequences}
\LANGen{
\Question{
A sequence has \(n\)th term \(3n^2-4n+1\). In the sequence, there are two consecutive terms that sum to \(581\).  Find the bigger one of these two terms?}
{\(320\)}{\(11\)}{\(261\)}{\(10\)}{}{}}
\LANGpl{
\Question{
Ciąg ma \(n\)-ty wyraz \(3n^2-4n+1\). W ciągu występują dwa następujące po sobie wyrazy, które sumują się do \(581\). Znajdź większy z tych dwóch wyrazów?}
{\(320\)}{\(11\)}{\(261\)}{\(10\)}{}{}}
\LANGcs{
\Question{
Posloupnost je dána \(n\)-tým členem \(3n^2-4n+1\). Součet dvou po sobě jdoucích členů je \(581\).  Určete větší z nich.}
{\(320\)}{\(11\)}{\(261\)}{\(10\)}{}{}}
\LANGsk{
\Question{
Postupnosť je daná \(n\)-tým členom \(3n^2-4n+1\). Súčet dvoch po sebe idúcich  členov je \(581\). Určte väčší z nich.}
{\(320\)}{\(11\)}{\(261\)}{\(10\)}{}{}}
\LANGes{
\Question{
Una sucesión tiene \(n\) el término \(3n^2-4n+1\). En la sucesión hay dos términos consecutivos cuya suma es \(581\).  Halla el mayor de estos dos términos.}
{\(320\)}{\(11\)}{\(261\)}{\(10\)}{}{}}
\End

\Data\ID{2000010308}
\Updated{2022-08-25 10:08:26}
\Level{A}
\SubArea{Properties of Sequences}
\LANGen{
\Question{
A sequence has \(n\)th term \(\frac{1}{3+2n}\) . Which term of the sequence is the first to have the value less than \(\frac1{200}\)?}
{\( a_{99} \)}{\( a_{98} \)}{\( a_{101} \)}{\( a_{102}\)}{}{}}
\LANGpl{
\Question{
Ciąg ma \(n\)-ty wyraz \(\frac{1}{3+2n}\) . Który wyraz ciągu jako pierwszy ma wartość mniejszą niż \(\frac1{200}\)?}
{\( a_{99} \)}{\( a_{98} \)}{\( a_{101} \)}{\( a_{102}\)}{}{}}
\LANGcs{
\Question{
Který člen posloupnosti \( \left( \frac{1}{3+2n}\right)^{\infty}_{n=1} \) je jako první menší než \(\frac1{200}\)?}
{\( a_{99} \)}{\( a_{98} \)}{\( a_{101} \)}{\( a_{102}\)}{}{}}
\LANGsk{
\Question{
Počnúc ktorým členom sú všetky členy postupnosti \( \left( \frac{1}{3+2n}\right)^{\infty}_{n=1} \)  menšie ako \(\frac1{200}\)?}
{\( a_{99} \)}{\( a_{98} \)}{\( a_{101} \)}{\( a_{102}\)}{}{}}
\LANGes{
\Question{
Una sucesión tiene el siguiente término general: \(\frac{1}{3+2n}\) . ¿Qué término de la sucesión es el primero teniendo el valor menor que \(\frac1{200}\)?}
{\( a_{99} \)}{\( a_{98} \)}{\( a_{101} \)}{\( a_{102}\)}{}{}}
\End

\Data\ID{2000010309}
\Updated{2022-08-25 10:08:26}
\Level{A}
\SubArea{Properties of Sequences}
\LANGen{
\Question{
A sequence has \(n\)th term \(50-\frac{1}{2}n^2\). Which term of the sequence is the first to have the value less than \(0\)?}
{\( a_{11} \)}{\( a_{10} \)}{\( a_{6} \)}{\( a_{5}\)}{}{}}
\LANGpl{
\Question{
Ciąg ma \(n\)-ty wyraz \(50-\frac{1}{2}n^2\). Który wyraz ciągu jako pierwszy ma wartość mniejszą niż \(0\)?}
{\( a_{11} \)}{\( a_{10} \)}{\( a_{6} \)}{\( a_{5}\)}{}{}}
\LANGcs{
\Question{
Který člen posloupnosti \( \left(50-\frac{1}{2}n^2\right)^{\infty}_{n=1} \)  je jako první menší než \(0\)?}
{\( a_{11} \)}{\( a_{10} \)}{\( a_{6} \)}{\( a_{5}\)}{}{}}
\LANGsk{
\Question{
Počnúc ktorým členom sú všetky členy postupnosti \( \left(50-\frac{1}{2}n^2\right)^{\infty}_{n=1} \)  menšie ako \(0\)?}
{\( a_{11} \)}{\( a_{10} \)}{\( a_{6} \)}{\( a_{5}\)}{}{}}
\LANGes{
\Question{
Una sucesión tiene el siguiente término general: \(50-\frac{1}{2}n^2\). ¿Qué término de la sucesión es el primero teniendo el valor menor que \(0\)?}
{\( a_{11} \)}{\( a_{10} \)}{\( a_{6} \)}{\( a_{5}\)}{}{}}
\End

\Data\ID{2000010310}
\Updated{2022-08-25 10:08:26}
\Level{A}
\SubArea{Properties of Sequences}
\LANGen{
\Question{
The \(n\)th term of a sequence is \(\left(\frac12\right)^n\). Find the difference between the \(5\)th and the \(8\)th term of the sequence.}
{\( \frac{7}{256} \)}{\( \frac{3}{128} \)}{\(-\frac{7}{256} \)}{\( 0\)}{}{}}
\LANGpl{
\Question{
Ciąg ma \(n\)-ty wyraz \(\left(\frac12\right)^n\). Znajdź różnicę między piątym a ósmym wyrazem ciągu.}
{\( \frac{7}{256} \)}{\( \frac{3}{128} \)}{\(-\frac{7}{256} \)}{\( 0\)}{}{}}
\LANGcs{
\Question{
Posloupnost je určená vzorcem pro \(n\)-tý člen \(\left(\frac12\right)^n\). Určete rozdíl mezi \(5\). a \(8\). členem posloupnosti.}
{\( \frac{7}{256} \)}{\( \frac{3}{128} \)}{\(-\frac{7}{256} \)}{\( 0\)}{}{}}
\LANGsk{
\Question{
Postupnosť je daná \(n\)-tým členom \(\left(\frac12\right)^n\). Určte rozdiel medzi piatym a ôsmim členom tejto postupnosti.}
{\( \frac{7}{256} \)}{\( \frac{3}{128} \)}{\(-\frac{7}{256} \)}{\( 0\)}{}{}}
\LANGes{
\Question{
Una sucesión tiene el siguiente término general:  \(\left(\frac12\right)^n\). Halla la diferencia entre el quinto y el octavo término de la sucesión.}
{\( \frac{7}{256} \)}{\( \frac{3}{128} \)}{\(-\frac{7}{256} \)}{\( 0\)}{}{}}
\End

\Data\ID{2010000401}
\Updated{2022-08-25 10:08:59}
\Level{A}
\SubArea{Properties of Sequences}
\LANGen{
\Question{
We are given a sequence \( \left( \frac{n}{n+1} \right)_{n=1}^{\infty} \). 
Which of the following formulations describes how is the given sequence defined?}
{defined by a formula for the \(n\)th term}{defined by a list of the sequence elements}{defined by a graph of the sequence}{defined by a recursive formula for the sequence}{}{}}
\LANGpl{
\Question{
Otrzymujemy ciąg \( \left( \frac{n}{n+1} \right)_{n=1}^{\infty} \).
Które z poniższych sformułowań opisuje, jak zdefiniowany jest dany ciąg?}
{zdefiniowany wzorem na \(n\)-ty wyraz}{zdefiniowany przez listę elementów sekwencji}{zdefiniowany przez wykres sekwencji}{zdefiniowany przez rekurencyjną formułę dla sekwencji}{}{}}
\LANGcs{
\Question{
Je dána posloupnost \( \left( \frac{n}{n+1} \right)_{n=1}^{\infty} \). 
Vyberte možnost, která co nejlépe popisuje způsob zadání této posloupnosti.}
{vzorec pro \(n\)-tý člen}{výběr členů posloupnosti}{graf posloupnosti}{rekurentní určení posloupnosti}{}{}}
\LANGsk{
\Question{
Je daná postupnosť \( \left( \frac{n}{n+1} \right)_{n=1}^{\infty} \). 
Vyberte možnosť, ktorá čo najlepšie popisuje spôsob zadania tejto postupnosti.}
{vzorec pre \(n\)-tý člen}{výber členov postupnosti}{graf postupnosti}{rekurentné vyjadrenie postupnosti}{}{}}
\LANGes{
\Question{
Dada la sucesión \( \left( \frac{n}{n+1} \right)_{n=1}^{\infty} \). ¿Cuál de las siguientes opciones se usa para definir la sucesión dada?}
{la fórmula de \(n\) términos}{una lista de términos de sucesión}{un gráfico de sucesión}{una fórmula recursiva de una sucesión}{}{}}
\End

\Data\ID{2010000403}
\Updated{2022-08-25 10:08:59}
\Level{A}
\SubArea{Properties of Sequences}
\LANGen{
\Question{
We are given a sequence \( \left( 5n-3\right)^{\infty}_{n=1} \). 
What does this formula express?}
{a sequence of all natural numbers which after dividing by \(5\) give the remainder \(2\)}{a sequence of all natural numbers which are divisible by \(3\)}{a sequence of all natural numbers which are divisible by \(5\)}{a sequence of all natural numbers which after dividing by \(5\) give the remainder \(3\)}{}{}}
\LANGpl{
\Question{
Otrzymaliśmy ciąg \( \left( 5n-3\right)^{\infty}_{n=1} \).
Co wyraża ten wzór?}
{ciąg wszystkich liczb naturalnych, które po podzieleniu przez \(5\) dają resztę \(2\)}{ciąg wszystkich liczb naturalnych podzielnych przez \(3\)}{ciąg wszystkich liczb naturalnych podzielnych przez \(5\)}{ciąg wszystkich liczb naturalnych, które po podzieleniu przez \(5\) dają resztę \(3\)}{}{}}
\LANGcs{
\Question{
Je dána posloupnost \( \left( 5n-3\right)^{\infty}_{n=1} \). 
Tento zápis vyjadřuje:}
{posloupnost všech přirozených čísel, která při dělení \(5\) dávají zbytek \(2\)}{posloupnost všech přirozených čísel dělitelných \(3\)}{posloupnost všech přirozených čísel dělitelných \(5\)}{posloupnost všech přirozených čísel, která při dělení \(5\) dávají zbytek \(3\)}{}{}}
\LANGsk{
\Question{
Je daná postupnosť \( \left( 5n-3\right)^{\infty}_{n=1} \). 
Tento zápis vyjadruje:}
{postupnosť všetkých prirodzených čísel, ktoré pri delení \(5\) dávajú zvyšok \(2\)}{postupnosť všetkých prirodzených čísel deliteľných \(3\)}{postupnosť všetkých prirodzených čísel deliteľných \(5\)}{postupnosť všetkých prirodzených čísel, ktoré pri delení \(5\) dávajú zvyšok \(3\)}{}{}}
\LANGes{
\Question{
Dada la sucesión \( \left( 5n-3\right)^{\infty}_{n=1} \). 
¿Qué define esta fórmula?}
{la secuencia de todos los números naturales que después de dividirlos por \(5\) dan el resto \(2\)}{la secuencia de todos los números naturales que son divisibles por \(3\)}{la secuencia de todos los números naturales que son divisibles por \(5\)}{la secuencia de todos los números naturales que después de dividirlos por \(5\) dan el resto \(3\)}{}{}}
\End

\Data\ID{2010000404}
\Updated{2022-08-25 10:08:59}
\Level{A}
\SubArea{Properties of Sequences}
\IMG{2010000401-figure0_0.pdf}
\LANGen{
\Question{
Which sequence is defined by the given graph?}
{\( \left( a_n \right)^{5}_{n=1} = 3,\ \ 2,\ \ 1,\ \ 2,\ \ 3 \)}{\( \left( a_n \right)^{10}_{n=1} = 1,\ \ 3,\ \ 2,\ \ 2,\ \ 3,\ \ 1,\ \ 4,\ \ 2,\ \ 5,\ \ 3 \)}{\( \left( a_n \right)^{5}_{n=1} = 1,\ \ 2,\ \ 3,\ \ 4,\ \ 5 \)}{\( \left( a_n \right)^{5}_{n=1} = 1,\ \ 2,\ \ 2,\ \ 3,\ \ 3 \)}{}{}}
\LANGpl{
\Question{
Który ciąg  jest wyrażony przez dany wykres?}
{\( \left( a_n \right)^{5}_{n=1} = 3,\ \ 2,\ \ 1,\ \ 2,\ \ 3 \)}{\( \left( a_n \right)^{10}_{n=1} = 1,\ \ 3,\ \ 2,\ \ 2,\ \ 3,\ \ 1,\ \ 4,\ \ 2,\ \ 5,\ \ 3 \)}{\( \left( a_n \right)^{5}_{n=1} = 1,\ \ 2,\ \ 3,\ \ 4,\ \ 5 \)}{\( \left( a_n \right)^{5}_{n=1} = 1,\ \ 2,\ \ 2,\ \ 3,\ \ 3 \)}{}{}}
\LANGcs{
\Question{
Vyberte posloupnost, která je daná grafem.}
{\( \left( a_n \right)^{5}_{n=1} = 3,\ \ 2,\ \ 1,\ \ 2,\ \ 3 \)}{\( \left( a_n \right)^{10}_{n=1} = 1,\ \ 3,\ \ 2,\ \ 2,\ \ 3,\ \ 1,\ \ 4,\ \ 2,\ \ 5,\ \ 3 \)}{\( \left( a_n \right)^{5}_{n=1} = 1,\ \ 2,\ \ 3,\ \ 4,\ \ 5 \)}{\( \left( a_n \right)^{5}_{n=1} = 1,\ \ 2,\ \ 2,\ \ 3,\ \ 3 \)}{}{}}
\LANGsk{
\Question{
Vyberte postupnosť, ktorá je daná grafom.}
{\( \left( a_n \right)^{5}_{n=1} = 3,\ \ 2,\ \ 1,\ \ 2,\ \ 3 \)}{\( \left( a_n \right)^{10}_{n=1} = 1,\ \ 3,\ \ 2,\ \ 2,\ \ 3,\ \ 1,\ \ 4,\ \ 2,\ \ 5,\ \ 3 \)}{\( \left( a_n \right)^{5}_{n=1} = 1,\ \ 2,\ \ 3,\ \ 4,\ \ 5 \)}{\( \left( a_n \right)^{5}_{n=1} = 1,\ \ 2,\ \ 2,\ \ 3,\ \ 3 \)}{}{}}
\LANGes{
\Question{
Elige la sucesión representada por el siguiente gráfico.}
{\( \left( a_n \right)^{5}_{n=1} = 3,\ \ 2,\ \ 1,\ \ 2,\ \ 3 \)}{\( \left( a_n \right)^{10}_{n=1} = 1,\ \ 3,\ \ 2,\ \ 2,\ \ 3,\ \ 1,\ \ 4,\ \ 2,\ \ 5,\ \ 3 \)}{\( \left( a_n \right)^{5}_{n=1} = 1,\ \ 2,\ \ 3,\ \ 4,\ \ 5 \)}{\( \left( a_n \right)^{5}_{n=1} = 1,\ \ 2,\ \ 2,\ \ 3,\ \ 3 \)}{}{}}
\End

\Data\ID{2010000405}
\Updated{2023-03-27 08:03:53}
\Level{A}
\SubArea{Properties of Sequences}
\IMG{2010000401-figure1_1.pdf}
\LANGen{
\Question{
We are given a sequence \( \left( a_n \right)^{5}_{n=1}\) defined by the following graph. Find the formula of the \(n\)th term of this sequence.}
{\( a_n = 3-2n,\ n \in\{1,\ \ 2,\ \ 3,\ \ 4,\ \ 5 \}\)}{\( a_n = 2n,\ n\in\{1,\ \ 2,\ \ 3,\ \ 4,\ \ 5 \}\)}{\( a_n = 1-2n,\ n\in\{1,\ \ 2,\ \ 3,\ \ 4,\ \ 5 \}\)}{\( a_n = 2n-3,\ n\in\{1,\ \ 2,\ \ 3,\ \ 4,\ \ 5 \}\)}{}{}}
\LANGpl{
\Question{
Dany jest ciąg \( \left( a_n \right)^{5}_{n=1}\) zdefiniowany na poniższym wykresie. Znajdź wzór \(n\)-tego wyrazu tego ciągu.}
{\( a_n = 3-2n,\ n \in\{1,\ \ 2,\ \ 3,\ \ 4,\ \ 5 \}\)}{\( a_n = 2n,\ n\in\{1,\ \ 2,\ \ 3,\ \ 4,\ \ 5 \}\)}{\( a_n = 1-2n,\ n\in\{1,\ \ 2,\ \ 3,\ \ 4,\ \ 5 \}\)}{\( a_n = 2n-3,\ n\in\{1,\ \ 2,\ \ 3,\ \ 4,\ \ 5 \}\)}{}{}}
\LANGcs{
\Question{
Posloupnost \( \left( a_n \right)^{5}_{n=1}\) je definována grafem na obrázku. Vzorec pro \(n\)-tý člen této posloupnosti je:}
{\( a_n = 3-2n,\ n \in\{1,\ \ 2,\ \ 3,\ \ 4,\ \ 5 \}\)}{\( a_n = 2n,\ n\in\{1,\ \ 2,\ \ 3,\ \ 4,\ \ 5 \}\)}{\( a_n = 1-2n,\ n\in\{1,\ \ 2,\ \ 3,\ \ 4,\ \ 5 \}\)}{\( a_n = 2n-3,\ n\in\{1,\ \ 2,\ \ 3,\ \ 4,\ \ 5 \}\)}{}{}}
\LANGsk{
\Question{
Je daná postupnosť \( \left( a_n \right)^{5}_{n=1}\), ktorá je definovaná nasledujúcim grafom.  Vzorec pre \(n\)-tý člen tejto postupnosti je:}
{\( a_n = 3-2n,\ n \in\{1,\ \ 2,\ \ 3,\ \ 4,\ \ 5 \}\)}{\( a_n = 2n,\ n\in\{1,\ \ 2,\ \ 3,\ \ 4,\ \ 5 \}\)}{\( a_n = 1-2n,\ n\in\{1,\ \ 2,\ \ 3,\ \ 4,\ \ 5 \}\)}{\( a_n = 2n-3,\ n\in\{1,\ \ 2,\ \ 3,\ \ 4,\ \ 5 \}\)}{}{}}
\LANGes{
\Question{
Dada la sucesión \( \left( a_n \right)^{5}_{n=1}\) representada por el siguiente gráfico. Halla el término general de la sucesión.}
{\( a_n = 3-2n,\ n \in\{1,\ \ 2,\ \ 3,\ \ 4,\ \ 5 \}\)}{\( a_n = 2n,\ n\in\{1,\ \ 2,\ \ 3,\ \ 4,\ \ 5 \}\)}{\( a_n = 1-2n,\ n\in\{1,\ \ 2,\ \ 3,\ \ 4,\ \ 5 \}\)}{\( a_n = 2n-3,\ n\in\{1,\ \ 2,\ \ 3,\ \ 4,\ \ 5 \}\)}{}{}}
\End

\Data\ID{2010000406}
\Updated{2023-03-27 08:03:19}
\Level{A}
\SubArea{Properties of Sequences}
\IMG{2010000401-figure2_1.pdf}
\LANGen{
\Question{
We are given a sequence \( \left( a_n \right)^{5}_{n=1}\) defined by the following graph. Find the formula of the \(n\)th term of this sequence.}
{\( a_n = 2n-3,\ n \in\{1,\ \ 2,\ \ 3,\ \ 4,\ \ 5 \}\)}{\( a_n = 2n,\ n \in\{1,\ \ 2,\ \ 3,\ \ 4,\ \ 5 \}\)}{\( a_n = 3-2n,\ n \in\{1,\ \ 2,\ \ 3,\ \ 4,\ \ 5 \}\)}{\( a_n = 2n-1,\ n \in\{1,\ \ 2,\ \ 3,\ \ 4,\ \ 5 \}\)}{}{}}
\LANGpl{
\Question{
Dany jest ciąg \( \left( a_n \right)^{5}_{n=1}\) zdefiniowany na poniższym wykresie. Znajdź wzór \(n\)-tego wyrazu tego ciągu.}
{\( a_n = 2n-3,\ n \in\{1,\ \ 2,\ \ 3,\ \ 4,\ \ 5 \}\)}{\( a_n = 2n,\ n \in\{1,\ \ 2,\ \ 3,\ \ 4,\ \ 5 \}\)}{\( a_n = 3-2n,\ n \in\{1,\ \ 2,\ \ 3,\ \ 4,\ \ 5 \}\)}{\( a_n = 2n-1,\ n \in\{1,\ \ 2,\ \ 3,\ \ 4,\ \ 5 \}\)}{}{}}
\LANGcs{
\Question{
Posloupnost \( \left( a_n \right)^{5}_{n=1}\) je definována grafem na obrázku. Vzorec pro \(n\)-tý člen této posloupnosti je:}
{\( a_n = 2n-3,\ n \in\{1,\ \ 2,\ \ 3,\ \ 4,\ \ 5 \}\)}{\( a_n = 2n,\ n \in\{1,\ \ 2,\ \ 3,\ \ 4,\ \ 5 \}\)}{\( a_n = 3-2n,\ n \in\{1,\ \ 2,\ \ 3,\ \ 4,\ \ 5 \}\)}{\( a_n = 2n-1,\ n \in\{1,\ \ 2,\ \ 3,\ \ 4,\ \ 5 \}\)}{}{}}
\LANGsk{
\Question{
Je daná postupnosť \( \left( a_n \right)^{5}_{n=1}\), ktorá je definovaná nasledujúcim grafom. Vzorec pre \(n\)-tý člen tejto postupnosti je:}
{\( a_n = 2n-3,\ n \in\{1,\ \ 2,\ \ 3,\ \ 4,\ \ 5 \}\)}{\( a_n = 2n,\ n \in\{1,\ \ 2,\ \ 3,\ \ 4,\ \ 5 \}\)}{\( a_n = 3-2n,\ n \in\{1,\ \ 2,\ \ 3,\ \ 4,\ \ 5 \}\)}{\( a_n = 2n-1,\ n \in\{1,\ \ 2,\ \ 3,\ \ 4,\ \ 5 \}\)}{}{}}
\LANGes{
\Question{
Dada la sucesión \( \left( a_n \right)^{5}_{n=1}\) representada por el siguiente gráfico. Halla el término general de la sucesión.}
{\( a_n = 2n-3,\ n \in\{1,\ \ 2,\ \ 3,\ \ 4,\ \ 5 \}\)}{\( a_n = 2n,\ n \in\{1,\ \ 2,\ \ 3,\ \ 4,\ \ 5 \}\)}{\( a_n = 3-2n,\ n \in\{1,\ \ 2,\ \ 3,\ \ 4,\ \ 5 \}\)}{\( a_n = 2n-1,\ n \in\{1,\ \ 2,\ \ 3,\ \ 4,\ \ 5 \}\)}{}{}}
\End

\Data\ID{2100005808}
\Updated{2022-07-20 07:07:27}
\Level{A}
\SubArea{Properties of Sequences}
\LANGen{
\Question{
In one of the pictures given below is \textbf{not} a graph of a sequence. Choose this graph.}
{}{}{}{}{}{}}
\LANGpl{
\Question{
Jeden z poniższych obrazków \textbf{nie} jest wykresem funkcji. Wybierz ten wykres.}
{}{}{}{}{}{}}
\LANGcs{
\Question{
Na kterém z obrázků \textbf{není} graf posloupnosti?}
{}{}{}{}{}{}}
\LANGsk{
\Question{
Na ktorom z nasledujúcich obrázkov \textbf{nie je} graf postupnosti?}
{}{}{}{}{}{}}
\LANGes{
\Question{
Una de las siguientes imágenes \textbf{no} representa el gráfico de una sucesión. ¿Cuál?}
{}{}{}{}{}{}}

\IMGanswer{2010005801-figure2.pdf}
\IMGanswer{2010005801-figure3.pdf}
\IMGanswer{2010005801-figure4.pdf}
\IMGanswer{2010005801-figure5_0.pdf}\End

\Data\ID{9000063803}
\Updated{2020-07-15 03:07:37}
\Level{A}
\SubArea{Properties of Sequences}
\LANGen{
\Question{
We are given the sequence \(\left (\cos n \frac{\pi }{4}\right )_{n=1}^{\infty }\).
Find the sum of the first six terms of this sequence.}
{\(-\frac{2+\sqrt{2}} 
 {2} \)}{\(-\frac{\sqrt{2}} 
 {2} \)}{\(- 1\)}{\(0\)}{}{}}
\LANGpl{
\Question{
Dany jest ciąg \(\left (\cos n \frac{\pi }{4}\right )_{n=1}^{\infty }\).
Oblicz sumę sześciu początkowych wyrazów tego ciągu.}
{\(-\frac{2+\sqrt{2}} 
 {2} \)}{\(-\frac{\sqrt{2}} 
 {2} \)}{\(- 1\)}{\(0\)}{}{}}
\LANGcs{
\Question{
Je dána posloupnost \(\left (\cos n \frac{\pi }{4}\right )_{n=1}^{\infty }\).
Součet prvních šesti členů této posloupnosti je roven:}
{\(-\frac{2+\sqrt{2}} 
 {2} \)}{\(-\frac{\sqrt{2}} 
 {2} \)}{\(- 1\)}{\(0\)}{}{}}
\LANGsk{
\Question{
Je daná postupnosť \(\left (\cos n \frac{\pi }{4}\right )_{n=1}^{\infty }\).
Súčet prvých šiestich členov tejto postupnosti je rovný:}
{\(-\frac{2+\sqrt{2}} 
 {2} \)}{\(-\frac{\sqrt{2}} 
 {2} \)}{\(- 1\)}{\(0\)}{}{}}
\LANGes{
\Question{
Dada la sucesión \(\left (\cos n \frac{\pi }{4}\right )_{n=1}^{\infty }\).
Halla la suma de los seis primeros términos de la sucesión.}
{\(-\frac{2+\sqrt{2}} 
 {2} \)}{\(-\frac{\sqrt{2}} 
 {2} \)}{\(- 1\)}{\(0\)}{}{}}
\End

\Data\ID{9000063804}
\Updated{2020-07-15 03:07:37}
\Level{A}
\SubArea{Properties of Sequences}
\LANGen{
\Question{
We are given the sequence \(\left (\log 10^{n}\right )_{n=1}^{\infty }\).
Find the product of the first five terms of this sequence.}
{\(120\)}{\(0\)}{\(5\)}{\(6\)}{}{}}
\LANGpl{
\Question{
Dany jest ciąg  \(\left (\log 10^{n}\right )_{n=1}^{\infty }\).
Oblicz produkt pięciu początkowych wyrazów tego ciągu.}
{\(120\)}{\(0\)}{\(5\)}{\(6\)}{}{}}
\LANGcs{
\Question{
Je dána posloupnost \(\left (\log 10^{n}\right )_{n=1}^{\infty }\).
Součin prvních pěti členů této posloupnosti je roven:}
{\(120\)}{\(0\)}{\(5\)}{\(6\)}{}{}}
\LANGsk{
\Question{
Je daná postupnosť \(\left (\log 10^{n}\right )_{n=1}^{\infty }\).
Súčin prvých piatich členov tejto postupnosti je rovný:}
{\(120\)}{\(0\)}{\(5\)}{\(6\)}{}{}}
\LANGes{
\Question{
Dada la sucesión \(\left (\log 10^{n}\right )_{n=1}^{\infty }\).
Halla el producto de los cinco primeros términos de la sucesión.}
{\(120\)}{\(0\)}{\(5\)}{\(6\)}{}{}}
\End

\Data\ID{9000063805}
\Updated{2023-03-27 08:03:21}
\Level{A}
\SubArea{Properties of Sequences}
\LANGen{
\Question{
Consider sequence \(a_{n+1} = 2a_{n} - a_{n-1}\) 
with \(a_{1} = 3\) 
and \(a_{2} = 5\).
Find \(a_{3} + a_{4}\).}
{\(16\)}{\(12\)}{\(0\)}{\(- 2\)}{}{}}
\LANGpl{
\Question{
Rozważ ciąg  \(a_{n+1} = 2a_{n} - a_{n-1}\) 
z \(a_{1} = 3\) 
i \(a_{2} = 5\).
Oblicz \(a_{3} + a_{4}\).}
{\(16\)}{\(12\)}{\(0\)}{\(- 2\)}{}{}}
\LANGcs{
\Question{
Uvažujme rekurentně zadanou posloupnost
\(a_{n+1} = 2a_{n} - a_{n-1}\), kde
\(a_{1} = 3\) a
\(a_{2} = 5\).
Potom platí:}
{\(a_{3} + a_{4} = 16\)}{\(a_{3} + a_{4} = 12\)}{\(a_{3} + a_{4} = 0\)}{\(a_{3} + a_{4} = -2\)}{}{}}
\LANGsk{
\Question{
Je daná rekurentne zadaná postupnosť
\(a_{n+1} = 2a_{n} - a_{n-1}\), kde
\(a_{1} = 3\) a
\(a_{2} = 5\).
Potom platí:}
{\(a_{3} + a_{4} = 16\)}{\(a_{3} + a_{4} = 12\)}{\(a_{3} + a_{4} = 0\)}{\(a_{3} + a_{4} = -2\)}{}{}}
\LANGes{
\Question{
Dada la sucesión \(a_{n+1} = 2a_{n} - a_{n-1}\) 
con \(a_{1} = 3\) 
y \(a_{2} = 5\).
Halla \(a_{3} + a_{4}\).}
{\(16\)}{\(12\)}{\(0\)}{\(- 2\)}{}{}}
\End

\Data\ID{9000063807}
\Updated{2020-07-15 03:07:37}
\Level{A}
\SubArea{Properties of Sequences}
\LANGen{
\Question{
Among the numbers \(5\),
\(15\),
\(28\),
\(47\) 
identify the number which is not a member of the sequence
\(\left (2n^{2} - 3\right )_{n=1}^{\infty }\).}
{\(28\)}{\(5\)}{\(15\)}{\(47\)}{}{}}
\LANGpl{
\Question{
Spośród liczb \(5\),
\(15\),
\(28\),
\(47\) 
wybierz liczbę, która nie jest częścią tego ciągu
\(\left (2n^{2} - 3\right )_{n=1}^{\infty }\).}
{\(28\)}{\(5\)}{\(15\)}{\(47\)}{}{}}
\LANGcs{
\Question{
Které z čísel \(5\),
\(15\),
\(28\),
\(47\) není členem
posloupnosti \(\left (2n^{2} - 3\right )_{n=1}^{\infty }\)?}
{\(28\)}{\(5\)}{\(15\)}{\(47\)}{}{}}
\LANGsk{
\Question{
Ktoré z čísel \(5\),
\(15\),
\(28\),
\(47\) nie je členom
postupnosti \(\left (2n^{2} - 3\right )_{n=1}^{\infty }\)?}
{\(28\)}{\(5\)}{\(15\)}{\(47\)}{}{}}
\LANGes{
\Question{
¿Cuál de los siguientes números \(5\),
\(15\),
\(28\),
\(47\) 
no es uno de los términos de la sucesión dada?
\(\left (2n^{2} - 3\right )_{n=1}^{\infty }\).}
{\(28\)}{\(5\)}{\(15\)}{\(47\)}{}{}}
\End

\Data\ID{9000063810}
\Updated{2020-07-15 03:07:37}
\Level{A}
\SubArea{Properties of Sequences}
\LANGen{
\Question{
Consider the sequences \(\left (a_{n}\right )_{n=1}^{\infty }\) 
and \(\left (b_{n}\right )_{n=1}^{\infty }\) 
where \(a_{n} = 2^{n}\) 
and \(b_{n} = n^{2} - 1\),
respectively. Identify a true statement in the terms of these sequences.}
{\(a_{3} = b_{3}\)}{\(a_{2} = b_{2} + 2\)}{\(a_{4} = b_{4} - 2\)}{\(a_{5} = b_{5} - 8\)}{}{}}
\LANGpl{
\Question{
Rozważmy ciągi  \(\left (a_{n}\right )_{n=1}^{\infty }\) 
i \(\left (b_{n}\right )_{n=1}^{\infty }\) 
gdzie odpowiednio \(a_{n} = 2^{n}\) 
i \(b_{n} = n^{2} - 1\).
Która z poniższych odpowiedzi jest prawidłowa?}
{\(a_{3} = b_{3}\)}{\(a_{2} = b_{2} + 2\)}{\(a_{4} = b_{4} - 2\)}{\(a_{5} = b_{5} - 8\)}{}{}}
\LANGcs{
\Question{
Jsou dány posloupnosti \(\left (a_{n}\right )_{n=1}^{\infty }\),
kde \(a_{n} = 2^{n}\), a
\(\left (b_{n}\right )_{n=1}^{\infty }\), kde
\(b_{n} = n^{2} - 1\).
Potom platí:}
{\(a_{3} = b_{3}\)}{\(a_{2} = b_{2} + 2\)}{\(a_{4} = b_{4} - 2\)}{\(a_{5} = b_{5} - 8\)}{}{}}
\LANGsk{
\Question{
Sú dané postupnosti \(\left (a_{n}\right )_{n=1}^{\infty }\),
kde \(a_{n} = 2^{n}\), a
\(\left (b_{n}\right )_{n=1}^{\infty }\), kde
\(b_{n} = n^{2} - 1\).
Potom platí:}
{\(a_{3} = b_{3}\)}{\(a_{2} = b_{2} + 2\)}{\(a_{4} = b_{4} - 2\)}{\(a_{5} = b_{5} - 8\)}{}{}}
\LANGes{
\Question{
Dadas las sucesiones \(\left (a_{n}\right )_{n=1}^{\infty }\) 
y \(\left (b_{n}\right )_{n=1}^{\infty }\) 
donde \(a_{n} = 2^{n}\) 
y \(b_{n} = n^{2} - 1\), Identifica la proposición lógica con respecto a los términos de esta sucesión.}
{\(a_{3} = b_{3}\)}{\(a_{2} = b_{2} + 2\)}{\(a_{4} = b_{4} - 2\)}{\(a_{5} = b_{5} - 8\)}{}{}}
\End

\Data\ID{1003084909}
\Updated{2020-07-15 03:07:21}
\Level{B}
\SubArea{Properties of Sequences}
\LANGen{
\Question{
We were given an oscillating sequence \( 3\text{, }-3\text{, }\ 3\text{, }-3\text{, }\ 3\dots \) (numbers \( 3 \) and \( -3 \) alternate regularly). What is the formula for the $n$th element of the sequence?}
{\( a_n=(-1)^{n+1}\cdot3\text{, }\ n\in\mathbb{N} \)}{\( a_n=(-1)^{n}\cdot3\text{, }\ n\in\mathbb{N} \)}{\( a_n=3^n\text{, }\ n\in\mathbb{N} \)}{\( a_n=-3^n\text{, }\ n\in\mathbb{N} \)}{}{}}
\LANGpl{
\Question{
Dany jest ciąg oscylacyjny \( 3\text{, }-3\text{, }\ 3\text{, }-3\text{, }\ 3\dots \) (liczby \( 3 \) i \( -3 \) zmieniają się regularnie). Jaki jest wzór na $n$-ty wyraz tego ciągu?}
{\( a_n=(-1)^{n+1}\cdot3\text{, }\ n\in\mathbb{N} \)}{\( a_n=(-1)^{n}\cdot3\text{, }\ n\in\mathbb{N} \)}{\( a_n=3^n\text{, }\ n\in\mathbb{N} \)}{\( a_n=-3^n\text{, }\ n\in\mathbb{N} \)}{}{}}
\LANGcs{
\Question{
Je dána oscilující posloupnost \( 3\text{, }-3\text{, }\ 3\text{, }-3\text{, }\ 3\dots \) (čísla \( 3 \) a \( -3 \) se pravidelně střídají). Vzorec pro $n$-tý člen této posloupnosti je:}
{\( a_n=(-1)^{n+1}\cdot3\text{, }\ n\in\mathbb{N} \)}{\( a_n=(-1)^{n}\cdot3\text{, }\ n\in\mathbb{N} \)}{\( a_n=3^n\text{, }\ n\in\mathbb{N} \)}{\( a_n=-3^n\text{, }\ n\in\mathbb{N} \)}{}{}}
\LANGsk{
\Question{
Je daná oscilujúca postupnosť \( 3\text{, }-3\text{, }\ 3\text{, }-3\text{, }\ 3\dots \) (čísla \( 3 \) a \( -3 \) sa pravidelne striedajú). Vzorec pre $n$-tý člen tejto postupnosti je:}
{\( a_n=(-1)^{n+1}\cdot3\text{, }\ n\in\mathbb{N} \)}{\( a_n=(-1)^{n}\cdot3\text{, }\ n\in\mathbb{N} \)}{\( a_n=3^n\text{, }\ n\in\mathbb{N} \)}{\( a_n=-3^n\text{, }\ n\in\mathbb{N} \)}{}{}}
\LANGes{
\Question{
Dada la sucesión oscilante \( 3\text{, }-3\text{, }\ 3\text{, }-3\text{, }\ 3\dots \) (los números \( 3 \) y \( -3 \) alternan regularmente). Halla el término general de la sucesión.}
{\( a_n=(-1)^{n+1}\cdot3\text{, }\ n\in\mathbb{N} \)}{\( a_n=(-1)^{n}\cdot3\text{, }\ n\in\mathbb{N} \)}{\( a_n=3^n\text{, }\ n\in\mathbb{N} \)}{\( a_n=-3^n\text{, }\ n\in\mathbb{N} \)}{}{}}
\End

\Data\ID{1003107301}
\Updated{2022-04-23 05:04:02}
\Level{B}
\SubArea{Properties of Sequences}
\LANGen{
\Question{
We are given the sequence \( \left( \frac{n+1}n \right)^{\infty}_{n=1} \). Find the recursive formula of such sequence.}
{\( a_1=2\,;\ a_{n+1}=a_n\frac{n(n+2)}{(n+1)^2},\ n\in\mathbb{N} \)}{\( a_1=2\,;\ a_{n+1}=a_n\frac{n(n+2)}{(n+1)},\ n\in\mathbb{N} \)}{\( a_1=1\,;\ a_{n+1}=a_n\frac{n(n+2)}{(n+1)^2},\ n\in\mathbb{N} \)}{\( a_1=2\,;\ a_{n+1}=a_n\frac{n(n-2)}{(n+1)^2},\ n\in\mathbb{N} \)}{}{}}
\LANGpl{
\Question{
Dany jest ciąg \( \left( \frac{n+1}n \right)^{\infty}_{n=1} \). Określ rekurencyjny wzór tego ciągu.}
{\( a_1=2\,;\ a_{n+1}=a_n\frac{n(n+2)}{(n+1)^2},\ n\in\mathbb{N} \)}{\( a_1=2\,;\ a_{n+1}=a_n\frac{n(n+2)}{(n+1)},\ n\in\mathbb{N} \)}{\( a_1=1\,;\ a_{n+1}=a_n\frac{n(n+2)}{(n+1)^2},\ n\in\mathbb{N} \)}{\( a_1=2\,;\ a_{n+1}=a_n\frac{n(n-2)}{(n+1)^2},\ n\in\mathbb{N} \)}{}{}}
\LANGcs{
\Question{
Je dána posloupnost  \( \left( \frac{n+1}n \right)^{\infty}_{n=1} \). Rekurentní vyjádření této posloupnosti je:}
{\( a_1=2\,;\ a_{n+1}=a_n\frac{n(n+2)}{(n+1)^2},\ n\in\mathbb{N} \)}{\( a_1=2\,;\ a_{n+1}=a_n\frac{n(n+2)}{(n+1)},\ n\in\mathbb{N} \)}{\( a_1=1\,;\ a_{n+1}=a_n\frac{n(n+2)}{(n+1)^2},\ n\in\mathbb{N} \)}{\( a_1=2\,;\ a_{n+1}=a_n\frac{n(n-2)}{(n+1)^2},\ n\in\mathbb{N} \)}{}{}}
\LANGsk{
\Question{
Je daná postupnosť  \( \left( \frac{n+1}n \right)^{\infty}_{n=1} \). Rekurentné vyjadrenie tejto postupnosti je:}
{\( a_1=2\,;\ a_{n+1}=a_n\frac{n(n+2)}{(n+1)^2},\ n\in\mathbb{N} \)}{\( a_1=2\,;\ a_{n+1}=a_n\frac{n(n+2)}{(n+1)},\ n\in\mathbb{N} \)}{\( a_1=1\,;\ a_{n+1}=a_n\frac{n(n+2)}{(n+1)^2},\ n\in\mathbb{N} \)}{\( a_1=2\,;\ a_{n+1}=a_n\frac{n(n-2)}{(n+1)^2},\ n\in\mathbb{N} \)}{}{}}
\LANGes{
\Question{
Dada la sucesión \( \left( \frac{n+1}n \right)^{\infty}_{n=1} \). Halla la fórmula recursiva de la sucesión.}
{\( a_1=2\,;\ a_{n+1}=a_n\frac{n(n+2)}{(n+1)^2},\ n\in\mathbb{N} \)}{\( a_1=2\,;\ a_{n+1}=a_n\frac{n(n+2)}{(n+1)},\ n\in\mathbb{N} \)}{\( a_1=1\,;\ a_{n+1}=a_n\frac{n(n+2)}{(n+1)^2},\ n\in\mathbb{N} \)}{\( a_1=2\,;\ a_{n+1}=a_n\frac{n(n-2)}{(n+1)^2},\ n\in\mathbb{N} \)}{}{}}
\End

\Data\ID{1003107302}
\Updated{2020-07-15 03:07:12}
\Level{B}
\SubArea{Properties of Sequences}
\LANGen{
\Question{
We are given the  sequence  \( \left( \frac{\sqrt2}4\left( \sqrt2 -1 \right)^n \right)^{\infty}_{n=1} \). Find the recursive formula of such sequence.}
{\( a_1=\frac{\sqrt2}4\left(\sqrt2-1\right);\ a_{n+1}=a_n\left(\sqrt2-1\right),\ n\in\mathbb{N} \)}{\( a_1=\frac{\sqrt2}4;\ a_{n+1}=a_n\left(\sqrt2-1\right),\ n\in\mathbb{N} \)}{\( a_1=\sqrt2-1\,;\ a_{n+1}=a_n\left(\sqrt2-1\right),\ n\in\mathbb{N} \)}{\( a_1=\frac{\sqrt2}4\left(\sqrt2-1\right);\ a_{n+1}=a_n\left(\sqrt2-1\right)^2,\ n\in\mathbb{N} \)}{}{}}
\LANGpl{
\Question{
Dany jest ciąg \( \left( \frac{\sqrt2}4\left( \sqrt2 -1 \right)^n \right)^{\infty}_{n=1} \). Określ rekurencyjny wzór tego ciągu.}
{\( a_1=\frac{\sqrt2}4\left(\sqrt2-1\right);\ a_{n+1}=a_n\left(\sqrt2-1\right),\ n\in\mathbb{N} \)}{\( a_1=\frac{\sqrt2}4;\ a_{n+1}=a_n\left(\sqrt2-1\right),\ n\in\mathbb{N} \)}{\( a_1=\sqrt2-1\,;\ a_{n+1}=a_n\left(\sqrt2-1\right),\ n\in\mathbb{N} \)}{\( a_1=\frac{\sqrt2}4\left(\sqrt2-1\right);\ a_{n+1}=a_n\left(\sqrt2-1\right)^2,\ n\in\mathbb{N} \)}{}{}}
\LANGcs{
\Question{
Je dána posloupnost  \( \left( \frac{\sqrt2}4\left( \sqrt2 -1 \right)^n \right)^{\infty}_{n=1} \). Rekurentní vyjádření této posloupnosti je:}
{\( a_1=\frac{\sqrt2}4\left(\sqrt2-1\right);\ a_{n+1}=a_n\left(\sqrt2-1\right),\ n\in\mathbb{N} \)}{\( a_1=\frac{\sqrt2}4;\ a_{n+1}=a_n\left(\sqrt2-1\right),\ n\in\mathbb{N} \)}{\( a_1=\sqrt2-1\,;\ a_{n+1}=a_n\left(\sqrt2-1\right),\ n\in\mathbb{N} \)}{\( a_1=\frac{\sqrt2}4\left(\sqrt2-1\right);\ a_{n+1}=a_n\left(\sqrt2-1\right)^2,\ n\in\mathbb{N} \)}{}{}}
\LANGsk{
\Question{
Je daná postupnosť  \( \left( \frac{\sqrt2}4\left( \sqrt2 -1 \right)^n \right)^{\infty}_{n=1} \). Rekurentné vyjadrenie tejto postupnosti je:}
{\( a_1=\frac{\sqrt2}4\left(\sqrt2-1\right);\ a_{n+1}=a_n\left(\sqrt2-1\right),\ n\in\mathbb{N} \)}{\( a_1=\frac{\sqrt2}4;\ a_{n+1}=a_n\left(\sqrt2-1\right),\ n\in\mathbb{N} \)}{\( a_1=\sqrt2-1\,;\ a_{n+1}=a_n\left(\sqrt2-1\right),\ n\in\mathbb{N} \)}{\( a_1=\frac{\sqrt2}4\left(\sqrt2-1\right);\ a_{n+1}=a_n\left(\sqrt2-1\right)^2,\ n\in\mathbb{N} \)}{}{}}
\LANGes{
\Question{
Dada la sucesión \( \left( \frac{\sqrt2}4\left( \sqrt2 -1 \right)^n \right)^{\infty}_{n=1} \). Halla la fórmula recursiva de la sucesión.}
{\( a_1=\frac{\sqrt2}4\left(\sqrt2-1\right);\ a_{n+1}=a_n\left(\sqrt2-1\right),\ n\in\mathbb{N} \)}{\( a_1=\frac{\sqrt2}4;\ a_{n+1}=a_n\left(\sqrt2-1\right),\ n\in\mathbb{N} \)}{\( a_1=\sqrt2-1\,;\ a_{n+1}=a_n\left(\sqrt2-1\right),\ n\in\mathbb{N} \)}{\( a_1=\frac{\sqrt2}4\left(\sqrt2-1\right);\ a_{n+1}=a_n\left(\sqrt2-1\right)^2,\ n\in\mathbb{N} \)}{}{}}
\End

\Data\ID{1003107303}
\Updated{2022-04-23 05:04:57}
\Level{B}
\SubArea{Properties of Sequences}
\LANGen{
\Question{
We are given the  sequence \( \left( 2^{2-n} \right)^{\infty}_{n=1} \). Find the recursive formula of such sequence.}
{\( a_1=2\,;\ a_{n+1}=a_n\cdot\frac12,\ n\in\mathbb{N} \)}{\( a_1=2\,;\ a_{n+1}=a_n\cdot2,\ n\in\mathbb{N} \)}{\( a_1=2\,;\ a_{n+1}=a_n\cdot2^n,\ n\in\mathbb{N} \)}{\( a_1=2\,;\ a_{n+1}=a_n\cdot\frac1{2^n},\ n\in\mathbb{N} \)}{}{}}
\LANGpl{
\Question{
Dany jest ciąg \( \left( 2^{2-n} \right)^{\infty}_{n=1} \). Określ rekurencyjny wzór tego ciągu.}
{\( a_1=2\,;\ a_{n+1}=a_n\cdot\frac12,\ n\in\mathbb{N} \)}{\( a_1=2\,;\ a_{n+1}=a_n\cdot2,\ n\in\mathbb{N} \)}{\( a_1=2\,;\ a_{n+1}=a_n\cdot2^n,\ n\in\mathbb{N} \)}{\( a_1=2\,;\ a_{n+1}=a_n\cdot\frac1{2^n},\ n\in\mathbb{N} \)}{}{}}
\LANGcs{
\Question{
Je dána posloupnost \( \left( 2^{2-n} \right)^{\infty}_{n=1} \). Rekurentní vyjádření této posloupnosti je:}
{\( a_1=2\,;\ a_{n+1}=a_n\cdot\frac12,\ n\in\mathbb{N} \)}{\( a_1=2\,;\ a_{n+1}=a_n\cdot2,\ n\in\mathbb{N} \)}{\( a_1=2\,;\ a_{n+1}=a_n\cdot2^n,\ n\in\mathbb{N} \)}{\( a_1=2\,;\ a_{n+1}=a_n\cdot\frac1{2^n},\ n\in\mathbb{N} \)}{}{}}
\LANGsk{
\Question{
Je daná postupnosť \( \left( 2^{2-n} \right)^{\infty}_{n=1} \). Rekurentné vyjadrenie tejto postupnosti je:}
{\( a_1=2\,;\ a_{n+1}=a_n\cdot\frac12,\ n\in\mathbb{N} \)}{\( a_1=2\,;\ a_{n+1}=a_n\cdot2,\ n\in\mathbb{N} \)}{\( a_1=2\,;\ a_{n+1}=a_n\cdot2^n,\ n\in\mathbb{N} \)}{\( a_1=2\,;\ a_{n+1}=a_n\cdot\frac1{2^n},\ n\in\mathbb{N} \)}{}{}}
\LANGes{
\Question{
Dada la sucesión \( \left( 2^{2-n} \right)^{\infty}_{n=1} \). Halla la fórmula recursiva de la sucesión.}
{\( a_1=2\,;\ a_{n+1}=a_n\cdot\frac12,\ n\in\mathbb{N} \)}{\( a_1=2\,;\ a_{n+1}=a_n\cdot2,\ n\in\mathbb{N} \)}{\( a_1=2\,;\ a_{n+1}=a_n\cdot2^n,\ n\in\mathbb{N} \)}{\( a_1=2\,;\ a_{n+1}=a_n\cdot\frac1{2^n},\ n\in\mathbb{N} \)}{}{}}
\End

\Data\ID{1003107304}
\Updated{2020-07-15 03:07:12}
\Level{B}
\SubArea{Properties of Sequences}
\LANGen{
\Question{
We are given a sequence \( \left( a_n \right)^{\infty}_{n=1} \)  defined recursively by: \(a_1=0\,;\ a_{n+1}=2-a_n,\ n\in\mathbb{N} \).
Find the \( n \)th term of this  sequence.}
{\( a_n=1+(-1)^n,\ n\in\mathbb{N} \)}{\( a_n=1+(-1)^{n+1},\ n\in\mathbb{N} \)}{\( a_n=1+(-1)^{n-1},\ n\in\mathbb{N} \)}{\( a_n=1-1^n,\ n\in\mathbb{N} \)}{}{}}
\LANGpl{
\Question{
Dany jest ciąg  \( \left( a_n \right)^{\infty}_{n=1} \)  określony rekurencyjnie przez \(a_1=0\,;\ a_{n+1}=2-a_n,\ n\in\mathbb{N} \).
Wyznacz $n$-ty wyraz tego ciągu.}
{\( a_n=1+(-1)^n,\ n\in\mathbb{N} \)}{\( a_n=1+(-1)^{n+1},\ n\in\mathbb{N} \)}{\( a_n=1+(-1)^{n-1},\ n\in\mathbb{N} \)}{\( a_n=1-1^n,\ n\in\mathbb{N} \)}{}{}}
\LANGcs{
\Question{
Posloupnost  \( \left( a_n \right)^{\infty}_{n=1} \)  je určena rekurentně: \(a_1=0\,;\ a_{n+1}=2-a_n,\ n\in\mathbb{N} \).
Vzorec pro \( n \)-tý člen této posloupnosti je:}
{\( a_n=1+(-1)^n,\ n\in\mathbb{N} \)}{\( a_n=1+(-1)^{n+1},\ n\in\mathbb{N} \)}{\( a_n=1+(-1)^{n-1},\ n\in\mathbb{N} \)}{\( a_n=1-1^n,\ n\in\mathbb{N} \)}{}{}}
\LANGsk{
\Question{
Postupnosť  \( \left( a_n \right)^{\infty}_{n=1} \)  je určená rekurentne: \(a_1=0\,;\ a_{n+1}=2-a_n,\ n\in\mathbb{N} \).
Vzorec pre \( n \)-tý člen tejto postupnosti je:}
{\( a_n=1+(-1)^n,\ n\in\mathbb{N} \)}{\( a_n=1+(-1)^{n+1},\ n\in\mathbb{N} \)}{\( a_n=1+(-1)^{n-1},\ n\in\mathbb{N} \)}{\( a_n=1-1^n,\ n\in\mathbb{N} \)}{}{}}
\LANGes{
\Question{
La sucesión \( \left( a_n \right)^{\infty}_{n=1} \)  viene dada por la fórmula recursiva: \(a_1=0\,;\ a_{n+1}=2-a_n,\ n\in\mathbb{N} \).
Halla el término general de la sucesión.}
{\( a_n=1+(-1)^n,\ n\in\mathbb{N} \)}{\( a_n=1+(-1)^{n+1},\ n\in\mathbb{N} \)}{\( a_n=1+(-1)^{n-1},\ n\in\mathbb{N} \)}{\( a_n=1-1^n,\ n\in\mathbb{N} \)}{}{}}
\End

\Data\ID{1003107305}
\Updated{2022-04-23 05:04:57}
\Level{B}
\SubArea{Properties of Sequences}
\LANGen{
\Question{
We are given a sequence \( \left( a_n \right)^{\infty}_{n=1} \)  defined recursively by: \( a_1=5\,;\ a_{n+1}=a_n+4,\ n\in\mathbb{N} \).
Find the \( n \)th term of this  sequence.}
{\( a_n=4n+1,\ n\in\mathbb{N} \)}{\( a_n=4n-1,\ n\in\mathbb{N} \)}{\( a_n=4n,\ n\in\mathbb{N} \)}{\( a_n=5n,\ n\in\mathbb{N} \)}{}{}}
\LANGpl{
\Question{
Dany jest ciąg \( \left( a_n \right)^{\infty}_{n=1} \)  określony rekurencyjnie przez \( a_1=5\,;\ a_{n+1}=a_n+4,\ n\in\mathbb{N} \).
Wyznacz $n$-ty wyraz tego ciągu.}
{\( a_n=4n+1,\ n\in\mathbb{N} \)}{\( a_n=4n-1,\ n\in\mathbb{N} \)}{\( a_n=4n,\ n\in\mathbb{N} \)}{\( a_n=5n,\ n\in\mathbb{N} \)}{}{}}
\LANGcs{
\Question{
Posloupnost \( \left( a_n \right)^{\infty}_{n=1} \)  je určena rekurentně: \( a_1=5\,;\ a_{n+1}=a_n+4,\ n\in\mathbb{N} \).
Vzorec pro \( n \)-tý člen této posloupnosti je:}
{\( a_n=4n+1,\ n\in\mathbb{N} \)}{\( a_n=4n-1,\ n\in\mathbb{N} \)}{\( a_n=4n,\ n\in\mathbb{N} \)}{\( a_n=5n,\ n\in\mathbb{N} \)}{}{}}
\LANGsk{
\Question{
Postupnosť \( \left( a_n \right)^{\infty}_{n=1} \)  je určená rekurentne: \( a_1=5\,;\ a_{n+1}=a_n+4,\ n\in\mathbb{N} \).
Vzorec pre \( n \)-tý člen tejto postupnosti je:}
{\( a_n=4n+1,\ n\in\mathbb{N} \)}{\( a_n=4n-1,\ n\in\mathbb{N} \)}{\( a_n=4n,\ n\in\mathbb{N} \)}{\( a_n=5n,\ n\in\mathbb{N} \)}{}{}}
\LANGes{
\Question{
La sucesión \( \left( a_n \right)^{\infty}_{n=1} \)  viene dada por la fórmula recursiva: \( a_1=5\,;\ a_{n+1}=a_n+4,\ n\in\mathbb{N} \).
Halla el término general de la sucesión.}
{\( a_n=4n+1,\ n\in\mathbb{N} \)}{\( a_n=4n-1,\ n\in\mathbb{N} \)}{\( a_n=4n,\ n\in\mathbb{N} \)}{\( a_n=5n,\ n\in\mathbb{N} \)}{}{}}
\End

\Data\ID{1003107306}
\Updated{2020-09-30 08:09:48}
\Level{B}
\SubArea{Properties of Sequences}
\LANGen{
\Question{
We are given a sequence \( \left( a_n \right)^{\infty}_{n=1} \)  defined recursively by: \( a_1=1\,;\ a_{n+1}=2a_n,\ n\in\mathbb{N} \).
Find the \( n \)th term of this  sequence.}
{\( a_n=2^{n-1},\ n\in\mathbb{N} \)}{\( a_n=2^n,\ n\in\mathbb{N} \)}{\( a_n=2^{n+1},\ n\in\mathbb{N} \)}{\( a_n=2^n-1,\ n\in\mathbb{N} \)}{}{}}
\LANGpl{
\Question{
Dany jest ciąg \( \left( a_n \right)^{\infty}_{n=1} \)  określony rekurencyjnie przez \( a_1=1\,;\ a_{n+1}=2a_n,\ n\in\mathbb{N} \).
Wyznacz $n$-ty wyraz tego ciągu.}
{\( a_n=2^{n-1},\ n\in\mathbb{N} \)}{\( a_n=2^n,\ n\in\mathbb{N} \)}{\( a_n=2^{n+1},\ n\in\mathbb{N} \)}{\( a_n=2^n-1,\ n\in\mathbb{N} \)}{}{}}
\LANGcs{
\Question{
Posloupnost \( \left( a_n \right)^{\infty}_{n=1} \)  je určena rekurentně: \( a_1=1\,;\ a_{n+1}=2a_n,\ n\in\mathbb{N} \).
Vzorec pro \( n \)-tý člen této posloupnosti je:}
{\( a_n=2^{n-1},\ n\in\mathbb{N} \)}{\( a_n=2^n,\ n\in\mathbb{N} \)}{\( a_n=2^{n+1},\ n\in\mathbb{N} \)}{\( a_n=2^n-1,\ n\in\mathbb{N} \)}{}{}}
\LANGsk{
\Question{
Postupnosť \( \left( a_n \right)^{\infty}_{n=1} \)  je určená rekurentne: \( a_1=1\,;\ a_{n+1}=2a_n,\ n\in\mathbb{N} \).
Vzorec pre \( n \)-tý člen tejto postupnosti je:}
{\( a_n=2^{n-1},\ n\in\mathbb{N} \)}{\( a_n=2^n,\ n\in\mathbb{N} \)}{\( a_n=2^{n+1},\ n\in\mathbb{N} \)}{\( a_n=2^n-1,\ n\in\mathbb{N} \)}{}{}}
\LANGes{
\Question{
La sucesión \( \left( a_n \right)^{\infty}_{n=1} \)  viene dada por la fórmula recursiva: \( a_1=1\,;\ a_{n+1}=2a_n,\ n\in\mathbb{N} \).
Halla el término general de la sucesión.}
{\( a_n=2^{n-1},\ n\in\mathbb{N} \)}{\( a_n=2^n,\ n\in\mathbb{N} \)}{\( a_n=2^{n+1},\ n\in\mathbb{N} \)}{\( a_n=2^n-1,\ n\in\mathbb{N} \)}{}{}}
\End

\Data\ID{1003107307}
\Updated{2022-04-23 05:04:02}
\Level{B}
\SubArea{Properties of Sequences}
\LANGen{
\Question{
We are given a sequence \( \left( a_n \right)^{\infty}_{n=1} \)  defined recursively by: \( a_1=1,\ a_2=5\) and \(\ a_{n+2}=a_{n+1}-a_n+d\), where \(\ n\in\mathbb{N} \).

Find the value of an unknown constant  \( d\in\mathbb{R} \) and of the term \( a_5 \) if \( a_3 = 10 \).}
{\( d=6,\ a_5=7 \)}{\( d=6,\ a_5=6 \)}{\( d=7,\ a_5=6 \)}{\( d=7,\ a_5=7 \)}{}{}}
\LANGpl{
\Question{
Dany jest ciąg \( \left( a_n \right)^{\infty}_{n=1} \)  określony rekurencyjnie przez  \( a_1=1,\ a_2=5\,;\ a_{n+2}=a_{n+1}-a_n+d,\ n\in\mathbb{N} \).

Wyznacz wartość nieznanej stałej  \( d\in\mathbb{R} \) i wyrazu \( a_5 \),  jeśli \( a_3 = 10 \).}
{\( d=6,\ a_5=7 \)}{\( d=6,\ a_5=6 \)}{\( d=7,\ a_5=6 \)}{\( d=7,\ a_5=7 \)}{}{}}
\LANGcs{
\Question{
Posloupnost \( \left( a_n \right)^{\infty}_{n=1} \)  je určena rekurentně: \( a_1=1,\ a_2=5\,;\ a_{n+2}=a_{n+1}-a_n+d,\ n\in\mathbb{N} \).

Určete hodnotu neznámé konstanty  \( d\in\mathbb{R} \) a členu \( a_5 \), víte-li, že \( a_3 = 10 \).}
{\( d=6,\ a_5=7 \)}{\( d=6,\ a_5=6 \)}{\( d=7,\ a_5=6 \)}{\( d=7,\ a_5=7 \)}{}{}}
\LANGsk{
\Question{
Postupnosť \( \left( a_n \right)^{\infty}_{n=1} \)  je určená rekurentne: \( a_1=1,\ a_2=5\,;\ a_{n+2}=a_{n+1}-a_n+d,\ n\in\mathbb{N} \).

Určte hodnotu neznámej konštanty  \( d\in\mathbb{R} \) a člena \( a_5 \), ak viete, že \( a_3 = 10 \).}
{\( d=6,\ a_5=7 \)}{\( d=6,\ a_5=6 \)}{\( d=7,\ a_5=6 \)}{\( d=7,\ a_5=7 \)}{}{}}
\LANGes{
\Question{
La sucesión  \( \left( a_n \right)^{\infty}_{n=1} \)  viene dada por la fórmula recursiva: \( a_1=1,\ a_2=5\,;\ a_{n+2}=a_{n+1}-a_n+d,\ n\in\mathbb{N} \).

Halla el valor de una constante desconocida  \( d\in\mathbb{R} \) y del término \( a_5 \) suponiendo que \( a_3 = 10 \).}
{\( d=6,\ a_5=7 \)}{\( d=6,\ a_5=6 \)}{\( d=7,\ a_5=6 \)}{\( d=7,\ a_5=7 \)}{}{}}
\End

\Data\ID{1003107309}
\Updated{2020-09-30 08:09:35}
\Level{B}
\SubArea{Properties of Sequences}
\LANGen{
\Question{
We are given the  sequence \( \left(\log2^n \right)^{\infty}_{n=1} \).
Find the recursive formula of such sequence.}
{\( a_1=\log 2\,;\ a_{n+1}=a_n+\log 2,\ n\in\mathbb{N} \)}{\( a_1=\log 2\,;\ a_{n+1}=a_n\cdot\log 2,\ n\in\mathbb{N} \)}{\( a_1=\log 2\,;\ a_{n+1}=a_n-\log 2,\ n\in\mathbb{N} \)}{\( a_1=\log 2\,;\ a_{n+1}=a_n+\log 2^n,\ n\in\mathbb{N} \)}{}{}}
\LANGpl{
\Question{
Dany jest ciąg \( \left(\log2^n \right)^{\infty}_{n=1} \).
Określ rekurencyjny wzór tego ciągu.}
{\( a_1=\log 2\,;\ a_{n+1}=a_n+\log 2,\ n\in\mathbb{N} \)}{\( a_1=\log 2\,;\ a_{n+1}=a_n\cdot\log 2,\ n\in\mathbb{N} \)}{\( a_1=\log 2\,;\ a_{n+1}=a_n-\log 2,\ n\in\mathbb{N} \)}{\( a_1=\log 2\,;\ a_{n+1}=a_n+\log 2^n,\ n\in\mathbb{N} \)}{}{}}
\LANGcs{
\Question{
Je dána posloupnost \( \left(\log2^n \right)^{\infty}_{n=1} \).
Rekurentní vyjádření této posloupnosti je:}
{\( a_1=\log 2\,;\ a_{n+1}=a_n+\log 2,\ n\in\mathbb{N} \)}{\( a_1=\log 2\,;\ a_{n+1}=a_n\cdot\log 2,\ n\in\mathbb{N} \)}{\( a_1=\log 2\,;\ a_{n+1}=a_n-\log 2,\ n\in\mathbb{N} \)}{\( a_1=\log 2\,;\ a_{n+1}=a_n+\log 2^n,\ n\in\mathbb{N} \)}{}{}}
\LANGsk{
\Question{
Je daná postupnosť \( \left(\log2^n \right)^{\infty}_{n=1} \).
Rekurentné vyjadrenie tejto postupnosti je:}
{\( a_1=\log 2\,;\ a_{n+1}=a_n+\log 2,\ n\in\mathbb{N} \)}{\( a_1=\log 2\,;\ a_{n+1}=a_n\cdot\log 2,\ n\in\mathbb{N} \)}{\( a_1=\log 2\,;\ a_{n+1}=a_n-\log 2,\ n\in\mathbb{N} \)}{\( a_1=\log 2\,;\ a_{n+1}=a_n+\log 2^n,\ n\in\mathbb{N} \)}{}{}}
\LANGes{
\Question{
Dada la sucesión \( \left(\log2^n \right)^{\infty}_{n=1} \).
Halla la fórmula recursiva de la sucesión.}
{\( a_1=\log 2\,;\ a_{n+1}=a_n+\log 2,\ n\in\mathbb{N} \)}{\( a_1=\log 2\,;\ a_{n+1}=a_n\cdot\log 2,\ n\in\mathbb{N} \)}{\( a_1=\log 2\,;\ a_{n+1}=a_n-\log 2,\ n\in\mathbb{N} \)}{\( a_1=\log 2\,;\ a_{n+1}=a_n+\log 2^n,\ n\in\mathbb{N} \)}{}{}}
\End

\Data\ID{1003107401}
\Updated{2020-07-15 03:07:21}
\Level{B}
\SubArea{Properties of Sequences}
\LANGen{
\Question{
We were given a sequence \( \left(\frac{2n+1}{n+2}\right)_{n=1}^{\infty} \).
 This sequence is:}
{increasing}{decreasing}{non-increasing}{constant}{}{}}
\LANGpl{
\Question{
Dany jest ciąg \( \left(\frac{2n+1}{n+2}\right)_{n=1}^{\infty} \).
 Ten ciąg jest:}
{rosnący}{malejący}{nierosnący}{stały}{}{}}
\LANGcs{
\Question{
Je dána posloupnost \( \left(\frac{2n+1}{n+2}\right)_{n=1}^{\infty} \).
 Tato posloupnost je:}
{rostoucí}{klesající}{nerostoucí}{konstantní}{}{}}
\LANGsk{
\Question{
Je daná postupnosť \( \left(\frac{2n+1}{n+2}\right)_{n=1}^{\infty} \).
 Táto postupnosť je:}
{rastúca}{klesajúca}{nerastúca}{konštantná}{}{}}
\LANGes{
\Question{
Dada la sucesión \( \left(\frac{2n+1}{n+2}\right)_{n=1}^{\infty} \).
Se trata de una sucesión:}
{creciente}{decreciente}{no creciente}{constante}{}{}}
\End

\Data\ID{1003107402}
\Updated{2020-07-15 03:07:21}
\Level{B}
\SubArea{Properties of Sequences}
\LANGen{
\Question{
We were given a sequence \( \left(\frac1{n(n+1)}\right)_{n=1}^{\infty} \).
This sequence is:}
{decreasing}{non-decreasing}{increasing}{constant}{}{}}
\LANGpl{
\Question{
Dany jest ciąg \( \left(\frac1{n(n+1)}\right)_{n=1}^{\infty} \).
Ten ciąg jest:}
{malejący}{niemalejący}{rosnący}{stały}{}{}}
\LANGcs{
\Question{
Je dána posloupnost \( \left(\frac1{n(n+1)}\right)_{n=1}^{\infty} \).
Tato posloupnost je:}
{klesající}{neklesající}{rostoucí}{konstantní}{}{}}
\LANGsk{
\Question{
Je daná postupnosť \( \left(\frac1{n(n+1)}\right)_{n=1}^{\infty} \).
Táto postupnosť je:}
{klesajúca}{neklesajúca}{rastúca}{konštantná}{}{}}
\LANGes{
\Question{
Dada la sucesión \( \left(\frac1{n(n+1)}\right)_{n=1}^{\infty} \).
Se trata de una sucesión:}
{decreciente}{no creciente}{creciente}{constante}{}{}}
\End

\Data\ID{1003107403}
\Updated{2020-07-15 03:07:21}
\Level{B}
\SubArea{Properties of Sequences}
\LANGen{
\Question{
We were given a sequence \( \left(\frac{n\cdot x}{n+1}\right)_{n=1}^{\infty} \).
For what values of the parameter \( x \), \(x\in\mathbb{R}\), is the sequence decreasing?}
{\( x\in(-\infty; 0) \)}{\( x\in(-\infty;0] \)}{\( x\in(0;\infty) \)}{\( x\in[0;\infty) \)}{}{}}
\LANGpl{
\Question{
Dany jest ciąg \( \left(\frac{n\cdot x}{n+1}\right)_{n=1}^{\infty} \).
Dla jakich wartości parametru \( x \), \(x\in\mathbb{R}\), dany ciąg jest malejący?}
{\( x\in(-\infty; 0) \)}{\( x\in(-\infty;0\rangle \)}{\( x\in(0;\infty) \)}{\( x\in\langle0;\infty) \)}{}{}}
\LANGcs{
\Question{
Je dána posloupnost \( \left(\frac{n\cdot x}{n+1}\right)_{n=1}^{\infty} \).
Pro které hodnoty parametru  \(x\in\mathbb{R}\) je tato posloupnost klesající?}
{\( x\in(-\infty; 0) \)}{\( x\in(-\infty;0\rangle \)}{\( x\in(0;\infty) \)}{\( x\in\langle0;\infty) \)}{}{}}
\LANGsk{
\Question{
Je daná postupnosť \( \left(\frac{n\cdot x}{n+1}\right)_{n=1}^{\infty} \).
Pre ktoré hodnoty parametra  \(x\in\mathbb{R}\) je táto postupnosť klesajúca?}
{\( x\in(-\infty; 0) \)}{\( x\in(-\infty;0\rangle \)}{\( x\in(0;\infty) \)}{\( x\in\langle0;\infty) \)}{}{}}
\LANGes{
\Question{
Dada la sucesión \( \left(\frac{n\cdot x}{n+1}\right)_{n=1}^{\infty} \).
¿Para qué valores del parámetro \( x \), \(x\in\mathbb{R}\), es una sucesión decreciente?}
{\( x\in(-\infty; 0) \)}{\( x\in(-\infty;0] \)}{\( x\in(0;\infty) \)}{\( x\in[0;\infty) \)}{}{}}
\End

\Data\ID{1003107404}
\Updated{2020-07-15 03:07:21}
\Level{B}
\SubArea{Properties of Sequences}
\LANGen{
\Question{
We were given a sequence \( \left(2\cdot x^n\right)_{n=1}^{\infty} \).
For what values of the parameter \( x \), \( x\in\mathbb{R} \), is the sequence increasing?}
{\( x\in(1;\infty) \)}{\( x\in[1;\infty) \)}{\( x\in(-\infty;1) \)}{\( x\in(-\infty;1]\)}{}{}}
\LANGpl{
\Question{
Dany jest ciąg \( \left(2\cdot x^n\right)_{n=1}^{\infty} \).
Dla jakich wartości parametru  \( x \), \( x\in\mathbb{R} \), dany ciąg jest rosnący?}
{\( x\in(1;\infty) \)}{\( x\in\langle1;\infty) \)}{\( x\in(-\infty;1) \)}{\( x\in(-\infty;1\rangle\)}{}{}}
\LANGcs{
\Question{
Je dána posloupnost \( \left(2\cdot x^n\right)_{n=1}^{\infty} \).
Pro které hodnoty parametru \( x\in\mathbb{R} \) je tato posloupnost rostoucí?}
{\( x\in(1;\infty) \)}{\( x\in\langle1;\infty) \)}{\( x\in(-\infty;1) \)}{\( x\in(-\infty;1\rangle\)}{}{}}
\LANGsk{
\Question{
Je daná postupnosť \( \left(2\cdot x^n\right)_{n=1}^{\infty} \).
Pre ktoré hodnoty parametra \( x\in\mathbb{R} \) je táto postupnosť rastúca?}
{\( x\in(1;\infty) \)}{\( x\in\langle1;\infty) \)}{\( x\in(-\infty;1) \)}{\( x\in(-\infty;1\rangle\)}{}{}}
\LANGes{
\Question{
Dada la sucesión \( \left(2\cdot x^n\right)_{n=1}^{\infty} \).
¿Para qué valores del parámetro \( x \), \( x\in\mathbb{R} \), es una sucesión creciente?}
{\( x\in(1;\infty) \)}{\( x\in[1;\infty) \)}{\( x\in(-\infty;1) \)}{\( x\in(-\infty;1]\)}{}{}}
\End

\Data\ID{1003107405}
\Updated{2020-07-15 03:07:21}
\Level{B}
\SubArea{Properties of Sequences}
\LANGen{
\Question{
We were given a sequence \( \left( \log2^n\right)_{n=1}^{\infty} \).
 This sequence is:}
{increasing}{decreasing}{non-increasing}{bounded}{}{}}
\LANGpl{
\Question{
Dany jest ciąg \( \left( \log2^n\right)_{n=1}^{\infty} \).
 Ten ciąg jest:}
{rosnący}{malejący}{nierosnący}{ograniczony}{}{}}
\LANGcs{
\Question{
Je dána posloupnost \( \left( \log2^n\right)_{n=1}^{\infty} \).
 Tato posloupnost je:}
{rostoucí}{klesající}{nerostoucí}{omezená}{}{}}
\LANGsk{
\Question{
Je daná postupnosť \( \left( \log2^n\right)_{n=1}^{\infty} \).
 Táto postupnosť je:}
{rastúca}{klesajúca}{nerastúca}{ohraničená}{}{}}
\LANGes{
\Question{
Dada la sucesión \( \left( \log2^n\right)_{n=1}^{\infty} \). Se trata de una sucesión:}
{creciente}{decreciente}{no creciente}{acotada}{}{}}
\End

\Data\ID{1003107406}
\Updated{2020-07-15 03:07:21}
\Level{B}
\SubArea{Properties of Sequences}
\LANGen{
\Question{
We were given a sequence \( \left( (-1)^n\cdot n\right)_{n=1}^{\infty} \).
Complete the sentence: This sequence is ...}
{neither increasing nor decreasing.}{increasing.}{decreasing.}{lower bounded.}{}{}}
\LANGpl{
\Question{
Dany jest ciąg \( \left( (-1)^n\cdot n\right)_{n=1}^{\infty} \).
Uzupełnij zdanie: Ten ciąg jest ...}
{ani rosnący ani malejący.}{rosnący.}{malejący.}{ograniczony z dołu.}{}{}}
\LANGcs{
\Question{
Je dána posloupnost \( \left( (-1)^n\cdot n\right)_{n=1}^{\infty} \).
Dokončete větu: Tato posloupnost ...}
{není ani rostoucí, ani klesající.}{je rostoucí.}{je klesající.}{je zdola omezená.}{}{}}
\LANGsk{
\Question{
Je daná postupnosť \( \left( (-1)^n\cdot n\right)_{n=1}^{\infty} \).
Dokončite vetu: Táto postupnosť ...}
{nie je ani rastúca, ani klesajúca.}{je rastúca.}{je klesajúca.}{je zdola ohraničená.}{}{}}
\LANGes{
\Question{
Dada la sucesión \( \left( (-1)^n\cdot n\right)_{n=1}^{\infty} \).
Se trata de una sucesión:}
{ni creciente ni decreciente}{creciente}{decreciente}{acotada inferiormente}{}{}}
\End

\Data\ID{1003107407}
\Updated{2022-04-23 05:04:57}
\Level{B}
\SubArea{Properties of Sequences}
\LANGen{
\Question{
We were given a sequence  \( \left( \log n \right)_{n=1}^{\infty} \). This sequence is:}
{lower bounded}{upper bounded}{bounded}{decreasing}{}{}}
\LANGpl{
\Question{
Dany jest ciąg \( \left( \log n \right)_{n=1}^{\infty} \). Ten ciąg jest:}
{ograniczony z dołu}{ograniczony z góry}{ograniczony}{malejący}{}{}}
\LANGcs{
\Question{
Je dána posloupnost  \( \left( \log n \right)_{n=1}^{\infty} \). Tato posloupnost je:}
{zdola omezená}{shora omezená}{omezená}{klesající}{}{}}
\LANGsk{
\Question{
Je daná postupnosť  \( \left( \log n \right)_{n=1}^{\infty} \). Táto postupnosť je:}
{zdola ohraničená}{zhora ohraničená}{ohraničená}{klesajúca}{}{}}
\LANGes{
\Question{
Dada la sucesión \( \left( \log n \right)_{n=1}^{\infty} \). Se trata de una sucesión:}
{acotada inferiormente}{acotada superiormente}{acotada}{decreciente}{}{}}
\End

\Data\ID{1003107408}
\Updated{2020-07-15 03:07:21}
\Level{B}
\SubArea{Properties of Sequences}
\LANGen{
\Question{
A sequence \( \left(a_n\right)_{n=1}^{\infty} \) is defined recursively by: \( a_1=5;\ a_{n+1}=2a_n-1\text{, } n\in\mathbb{N} \).
This sequence is:}
{lower bounded}{upper bounded}{bounded}{decreasing}{}{}}
\LANGpl{
\Question{
Ciąg \( \left(a_n\right)_{n=1}^{\infty} \) jest określony rekurencyjnie: \( a_1=5;\ a_{n+1}=2a_n-1\text{, } n\in\mathbb{N} \).
Ten ciąg jest:}
{ograniczony z dołu}{ograniczony z góry}{ograniczony}{malejący}{}{}}
\LANGcs{
\Question{
Posloupnost \( \left(a_n\right)_{n=1}^{\infty} \) je určena rekurentně: \( a_1=5;\ a_{n+1}=2a_n-1\text{, } n\in\mathbb{N} \).
Tato posloupnost je:}
{zdola omezená}{shora omezená}{omezená}{klesající}{}{}}
\LANGsk{
\Question{
Postupnosť \( \left(a_n\right)_{n=1}^{\infty} \) je určená rekurentne: \( a_1=5;\ a_{n+1}=2a_n-1\text{, } n\in\mathbb{N} \).
Táto postupnosť je:}
{zdola ohraničená}{zhora ohraničená}{ohraničená}{klesajúca}{}{}}
\LANGes{
\Question{
La sucesión \( \left(a_n\right)_{n=1}^{\infty} \) viene dada por la fórmula recursiva: \( a_1=5;\ a_{n+1}=2a_n-1\text{, } n\in\mathbb{N} \).
Se trata de una sucesión:}
{acotada inferiormente}{acotada superiormente}{acotada}{decreciente}{}{}}
\End

\Data\ID{1003107409}
\Updated{2020-07-15 03:07:21}
\Level{B}
\SubArea{Properties of Sequences}
\LANGen{
\Question{
We were given a sequence \( \left( 3+\frac1{2n}\right)_{n=1}^{\infty} \).
 This sequence is:}
{bounded}{increasing}{constant}{non-decreasing}{}{}}
\LANGpl{
\Question{
Dany jest ciąg \( \left( 3+\frac1{2n}\right)_{n=1}^{\infty} \).
Ten ciąg jest:}
{ograniczony}{rosnący}{stały}{niemalejący}{}{}}
\LANGcs{
\Question{
Je dána posloupnost \( \left( 3+\frac1{2n}\right)_{n=1}^{\infty} \).
 Tato posloupnost je:}
{omezená}{rostoucí}{konstantní}{neklesající}{}{}}
\LANGsk{
\Question{
Je daná postupnosť \( \left( 3+\frac1{2n}\right)_{n=1}^{\infty} \).
 Táto postupnosť je:}
{ohraničená}{rastúca}{konštantná}{neklesajúca}{}{}}
\LANGes{
\Question{
Dada la sucesión \( \left( 3+\frac1{2n}\right)_{n=1}^{\infty} \).
Se trata de una sucesión:}
{acotada}{creciente}{constante}{no creciente}{}{}}
\End

\Data\ID{1003107410}
\Updated{2020-10-02 11:10:11}
\Level{B}
\SubArea{Properties of Sequences}
\LANGen{
\Question{
We were given a sequence \( \left(-n^2\right)_{n=1}^{\infty} \).
 This sequence is:}
{upper bounded}{lower bounded}{bounded}{increasing}{}{}}
\LANGpl{
\Question{
Dany jest ciąg \( \left(-n^2\right)_{n=1}^{\infty} \).
 Ten ciąg jest:}
{ograniczony z góry}{ograniczony z dołu}{ograniczony}{rosnący}{}{}}
\LANGcs{
\Question{
Je dána posloupnost \( \left(-n^2\right)_{n=1}^{\infty} \).
 Tato posloupnost je:}
{shora omezená}{zdola omezená}{omezená}{rostoucí}{}{}}
\LANGsk{
\Question{
Je daná postupnosť \( \left(-n^2\right)_{n=1}^{\infty} \).
 Táto postupnosť je:}
{zhora ohraničená}{zdola ohraničená}{ohraničená}{rastúca}{}{}}
\LANGes{
\Question{
Dada la sucesión \( \left(-n^2\right)_{n=1}^{\infty} \). Se trata de una sucesión:}
{acotada superiormente}{acotada inferiormente}{acotada}{creciente}{}{}}
\End

\Data\ID{2000010307}
\Updated{2023-03-27 11:03:38}
\Level{B}
\SubArea{Properties of Sequences}
\LANGen{
\Question{
Which of the following sequences given by the recursive formulas is not decreasing?}
{\( a_{n+1} = \frac{1}{a_n}\), \( a_1=5\)}{\( a_{n+1} = \sqrt{a_n}\), \( a_1=16\)}{\( a_{n+1} = 0.5\cdot {a_n}\), \( a_1=12\)}{\( a_{n+1} = \frac{a_n}{n}\), \( a_1=24\)}{}{}}
\LANGpl{
\Question{
Który z poniższych ciągów  podanych przez wzory rekurencyjne nie maleje?}
{\( a_{n+1} = \frac{1}{a_n}\), \( a_1=5\)}{\( a_{n+1} = \sqrt{a_n}\), \( a_1=16\)}{\( a_{n+1} = 0{,}5\cdot {a_n}\), \( a_1=12\)}{\( a_{n+1} = \frac{a_n}{n}\), \( a_1=24\)}{}{}}
\LANGcs{
\Question{
Která posloupnost určená rekurentně není klesající?}
{\( a_{n+1} = \frac{1}{a_n}\), \( a_1=5\)}{\( a_{n+1} = \sqrt{a_n}\), \( a_1=16\)}{\( a_{n+1} = 0{,}5\cdot {a_n}\), \( a_1=12\)}{\( a_{n+1} = \frac{a_n}{n}\), \( a_1=24\)}{}{}}
\LANGsk{
\Question{
Ktorá postupnosť daná rekurzívnym vzťahom nie je klesajúca?}
{\( a_{n+1} = \frac{1}{a_n}\), \( a_1=5\)}{\( a_{n+1} = \sqrt{a_n}\), \( a_1=16\)}{\( a_{n+1} = 0{,}5\cdot {a_n}\), \( a_1=12\)}{\( a_{n+1} = \frac{a_n}{n}\), \( a_1=24\)}{}{}}
\LANGes{
\Question{
¿Cuál de las siguientes sucesiones dadas por las fórmulas recursivas no es decreciente?}
{\( a_{n+1} = \frac{1}{a_n}\), \( a_1=5\)}{\( a_{n+1} = \sqrt{a_n}\), \( a_1=16\)}{\( a_{n+1} = 0.5\cdot {a_n}\), \( a_1=12\)}{\( a_{n+1} = \frac{a_n}{n}\), \( a_1=24\)}{}{}}
\End

\Data\ID{2010000402}
\Updated{2022-08-25 10:08:59}
\Level{B}
\SubArea{Properties of Sequences}
\LANGen{
\Question{
We are given the sequence \( \left( \frac{n}{n+1} \right)^{\infty}_{n=1} \). Find the recursive formula of such sequence.}
{\( a_1=\frac{1}{2}\,;\ a_{n+1}=a_n\frac{(n+1)^2}{n(n+2)},\ n\in\mathbb{N} \)}{\( a_1={2}\,;\ a_{n+1}=a_n\frac{(n+1)^2}{n(n+2)},\ n\in\mathbb{N} \)}{\( a_1=\frac{1}{2}\,;\ a_{n+1}=a_n\frac{n(n+1)}{(n+1)(n+2)},\ n\in\mathbb{N} \)}{\( a_1={2}\,;\ a_{n+1}=a_n\frac{n(n+1)}{(n+1)(n+2)},\ n\in\mathbb{N} \)}{}{}}
\LANGpl{
\Question{
Otrzymujemy ciąg \( \left( \frac{n}{n+1} \right)^{\infty}_{n=1} \). Znajdź wzór rekurencyjny takiego ciągu.}
{\( a_1=\frac{1}{2}\,;\ a_{n+1}=a_n\frac{(n+1)^2}{n(n+2)},\ n\in\mathbb{N} \)}{\( a_1={2}\,;\ a_{n+1}=a_n\frac{(n+1)^2}{n(n+2)},\ n\in\mathbb{N} \)}{\( a_1=\frac{1}{2}\,;\ a_{n+1}=a_n\frac{n(n+1)}{(n+1)(n+2)},\ n\in\mathbb{N} \)}{\( a_1={2}\,;\ a_{n+1}=a_n\frac{n(n+1)}{(n+1)(n+2)},\ n\in\mathbb{N} \)}{}{}}
\LANGcs{
\Question{
Je dána posloupnost \( \left( \frac{n}{n+1} \right)^{\infty}_{n=1} \). Rekurentní vyjádření této posloupnosti je:}
{\( a_1=\frac{1}{2}\,;\ a_{n+1}=a_n\frac{(n+1)^2}{n(n+2)},\ n\in\mathbb{N} \)}{\( a_1={2}\,;\ a_{n+1}=a_n\frac{(n+1)^2}{n(n+2)},\ n\in\mathbb{N} \)}{\( a_1=\frac{1}{2}\,;\ a_{n+1}=a_n\frac{n(n+1)}{(n+1)(n+2)},\ n\in\mathbb{N} \)}{\( a_1={2}\,;\ a_{n+1}=a_n\frac{n(n+1)}{(n+1)(n+2)},\ n\in\mathbb{N} \)}{}{}}
\LANGsk{
\Question{
Je daná postupnosť \( \left( \frac{n}{n+1} \right)^{\infty}_{n=1} \). Rekurentné vyjadrenie tejto postupnosti je:}
{\( a_1=\frac{1}{2}\,;\ a_{n+1}=a_n\frac{(n+1)^2}{n(n+2)},\ n\in\mathbb{N} \)}{\( a_1={2}\,;\ a_{n+1}=a_n\frac{(n+1)^2}{n(n+2)},\ n\in\mathbb{N} \)}{\( a_1=\frac{1}{2}\,;\ a_{n+1}=a_n\frac{n(n+1)}{(n+1)(n+2)},\ n\in\mathbb{N} \)}{\( a_1={2}\,;\ a_{n+1}=a_n\frac{n(n+1)}{(n+1)(n+2)},\ n\in\mathbb{N} \)}{}{}}
\LANGes{
\Question{
Dada la sucesión \( \left( \frac{n}{n+1} \right)^{\infty}_{n=1} \). Halla la fórmula recursiva de esta sucesión.}
{\( a_1=\frac{1}{2}\,;\ a_{n+1}=a_n\frac{(n+1)^2}{n(n+2)},\ n\in\mathbb{N} \)}{\( a_1={2}\,;\ a_{n+1}=a_n\frac{(n+1)^2}{n(n+2)},\ n\in\mathbb{N} \)}{\( a_1=\frac{1}{2}\,;\ a_{n+1}=a_n\frac{n(n+1)}{(n+1)(n+2)},\ n\in\mathbb{N} \)}{\( a_1={2}\,;\ a_{n+1}=a_n\frac{n(n+1)}{(n+1)(n+2)},\ n\in\mathbb{N} \)}{}{}}
\End

\Data\ID{2010000701}
\Updated{2022-08-25 10:08:02}
\Level{B}
\SubArea{Properties of Sequences}
\LANGen{
\Question{
We are given the sequence \(\left (an + b\right )_{n=1}^{\infty }\).
This sequence satisfies \(a_{7} - a_{2} = -10\). Use this information to find \(a\).}
{\(a = -2\)}{\(a = 2\)}{\(a = -1\)}{\(a = 1\)}{}{}}
\LANGpl{
\Question{
Dany jest ciąg \(\left (an + b\right )_{n=1}^{\infty }\).
Ta sekwencja spełnia \(a_{7} - a_{2} = -10\). Użyj tych informacji, aby znaleźć \(a\).}
{\(a = -2\)}{\(a = 2\)}{\(a = -1\)}{\(a = 1\)}{}{}}
\LANGcs{
\Question{
Je dána posloupnost \(\left (an + b\right )_{n=1}^{\infty }\),
ve které platí, že \(a_{7} - a_{2} = -10\). Určete \(a\).}
{\(a = -2\)}{\(a = 2\)}{\(a = -1\)}{\(a = 1\)}{}{}}
\LANGsk{
\Question{
Je daná postupnosť \(\left (an + b\right )_{n=1}^{\infty }\),
v ktorej platí, že \(a_{7} - a_{2} = -10\). Nájdite \(a\).}
{\(a = -2\)}{\(a = 2\)}{\(a = -1\)}{\(a = 1\)}{}{}}
\LANGes{
\Question{
Dada la sucesión \(\left (an + b\right )_{n=1}^{\infty }\)
en la que vale \(a_{7} - a_{2} = -10\). Halla \(a\).}
{\(a = -2\)}{\(a = 2\)}{\(a = -1\)}{\(a = 1\)}{}{}}
\End

\Data\ID{2010000702}
\Updated{2022-08-25 10:08:02}
\Level{B}
\SubArea{Properties of Sequences}
\LANGen{
\Question{
We are given a sequence \( \left( a_n \right)^{\infty}_{n=1} \)  defined recursively by: \( a_1=-1,\ a_2=0\) and \(\ a_{n+2}=a_{n}-a_{n+1}-d\), where \(\ n\in\mathbb{N} \).

Find the value of an unknown constant  \( d\in\mathbb{R} \) and of the term \( a_5 \) if \( a_3 = -4 \).}
{\( d=3,\ a_5=-8 \)}{\( d=5,\ a_5=-10 \)}{\( d=3,\ a_5=1\)}{\( d=5,\ a_5=-9 \)}{}{}}
\LANGpl{
\Question{
Dany jest ciąg \( \left( a_n \right)^{\infty}_{n=1} \) zdefiniowany rekurencyjnie przez: \( a_1=-1,\ a_2=0\) i \(\ a_{ n+2}=a_{n}-a_{n+1}-d\), gdzie \(\ n\in\mathbb{N} \).

Znajdź wartość nieznanej stałej \( d\in\mathbb{R} \) i wyrazu \( a_5 \) jeśli \( a_3 = -4 \).}
{\( d=3,\ a_5=-8 \)}{\( d=5,\ a_5=-10 \)}{\( d=3,\ a_5=1\)}{\( d=5,\ a_5=-9 \)}{}{}}
\LANGcs{
\Question{
Posloupnost \( \left( a_n \right)^{\infty}_{n=1} \) je určená rekurentně: \( a_1=-1,\ a_2=0\); \(\ a_{n+2}=a_{n}-a_{n+1}-d\), \(\ n\in\mathbb{N} \).

Určete hodnotu neznámé konstanty \( d\in\mathbb{R} \) a členu \( a_5 \), když víte, že \( a_3 = -4 \).}
{\( d=3,\ a_5=-8 \)}{\( d=5,\ a_5=-10 \)}{\( d=3,\ a_5=1\)}{\( d=5,\ a_5=-9 \)}{}{}}
\LANGsk{
\Question{
Postupnosť \( \left( a_n \right)^{\infty}_{n=1} \) je určená rekurentne: \( a_1=-1,\ a_2=0\); \(\ a_{n+2}=a_{n}-a_{n+1}-d\), \(\ n\in\mathbb{N} \).

Určte hodnotu neznámej konštanty \( d\in\mathbb{R} \) a člena \( a_5 \), ak viete, že \( a_3 = -4 \).}
{\( d=3,\ a_5=-8 \)}{\( d=5,\ a_5=-10 \)}{\( d=3,\ a_5=1\)}{\( d=5,\ a_5=-9 \)}{}{}}
\LANGes{
\Question{
La sucesión \( \left( a_n \right)^{\infty}_{n=1} \) viene dada por la fórmula recursiva: \( a_1=-1,\ a_2=0\) y \(\ a_{n+2}=a_{n}-a_{n+1}-d\), \(\ n\in\mathbb{N} \).
Halla el valor de una constante desconocida \( d\in\mathbb{R} \) y del término \( a_5 \) suponiendo que \( a_3 = -4 \).}
{\( d=3,\ a_5=-8 \)}{\( d=5,\ a_5=-10 \)}{\( d=3,\ a_5=1\)}{\( d=5,\ a_5=-9 \)}{}{}}
\End

\Data\ID{2010000703}
\Updated{2023-03-27 08:03:50}
\Level{B}
\SubArea{Properties of Sequences}
\LANGen{
\Question{
Consider the sequence \(2a_{n} = a_{n+1} - a_{n-1}\) 
with \(a_{3} = 2\) 
and \(a_{4} = 5\).
Find \(a_{2} - a_{1}\).}
{\( 1\)}{\(4\)}{\(-20\)}{\(-25\)}{}{}}
\LANGpl{
\Question{
Rozważmy ciąg \(2a_{n} = a_{n+1} - a_{n-1}\)
z \(a_{3} = 2\)
oraz \(a_{4} = 5\).
Znajdź \(a_{2} - a_{1}\).}
{\( 1\)}{\(4\)}{\(-20\)}{\(-25\)}{}{}}
\LANGcs{
\Question{
Uvažujme rekurentně zadanou posloupnost \(2a_{n} = a_{n+1} - a_{n-1}\), 
kde \(a_{3} = 2\) 
a \(a_{4} = 5\).
Potom platí:}
{\(a_{2} - a_{1} = 1\)}{\(a_{2} - a_{1} = 4\)}{\(a_{2} - a_{1} = -20\)}{\(a_{2} - a_{1} = -25\)}{}{}}
\LANGsk{
\Question{
Je daná rekurentne zadaná postupnosť \(2a_{n} = a_{n+1} - a_{n-1}\), 
kde \(a_{3} = 2\) 
a \(a_{4} = 5\).
Potom platí:}
{\(a_{2} - a_{1} = 1\)}{\(a_{2} - a_{1} = 4\)}{\(a_{2} - a_{1} = -20\)}{\(a_{2} - a_{1} = -25\)}{}{}}
\LANGes{
\Question{
Dada la sucesión \(2a_{n} = a_{n+1} - a_{n-1}\) 
con \(a_{3} = 2\) 
y \(a_{4} = 5\).
Halla \(a_{2} - a_{1}\).}
{\( 1\)}{\(4\)}{\(-20\)}{\(-25\)}{}{}}
\End

\Data\ID{2010000704}
\Updated{2022-08-25 10:08:02}
\Level{B}
\SubArea{Properties of Sequences}
\LANGen{
\Question{
The \(n\)th term of a sequence is given by \(a_n = \frac{n^2-1}{n+5}\). Which term of this sequence is equal to \(4\)?}
{the seventh term}{the third term}{the twenty-first term}{the fourth term}{}{}}
\LANGpl{
\Question{
Ciąg jest określony wzorem \(a_n = \frac{n^2-1}{n+5}\). Który wyraz tego ciągu jest równy \(4\)?}
{siódmy wyraz}{trzeci wyraz}{dwudziesty pierwszy wyraz}{czwarty wyraz}{}{}}
\LANGcs{
\Question{
Posloupnost je dána vzorcem pro \(n\)-tý člen ve tvaru \(a_n = \frac{n^2-1}{n+5}\). Který člen této posloupnosti je roven \(4\)?}
{sedmý}{třetí}{dvacátý první}{čtvrtý}{}{}}
\LANGsk{
\Question{
Postupnosť je daná vzorcom pre \(n\)-tý člen, v tvare \(a_n = \frac{n^2-1}{n+5}\). Ktorý člen tejto postupnosti je rovný \(4\)?}
{siedmy}{tretí}{dvadsiaty prvý}{štvrtý}{}{}}
\LANGes{
\Question{
El término general de una sucesión viene dado por \(a_n = \frac{n^2-1}{n+5}\). ¿Qué término de la sucesión equivale a \(4\)?}
{el séptimo término}{el tercer término}{el vigésimo primer término}{el cuarto término}{}{}}
\End

\Data\ID{2010000705}
\Updated{2022-08-25 10:08:02}
\Level{B}
\SubArea{Properties of Sequences}
\LANGen{
\Question{
The \(n\)th term of a sequence is given by \(a_n = \frac{n^2-4}{n+4}\). Which term of this sequence is equal to \(5\)?}
{the eighth term}{the third term}{the ninth term}{the fifth term}{}{}}
\LANGpl{
\Question{
Ciąg jest określony wzorem \(a_n = \frac{n^2-4}{n+4}\). Który wyraz tego ciągu jest równy \(5\)?}
{ósmy wyraz}{trzeci wyraz}{dziewiąty wyraz}{piąty wyraz}{}{}}
\LANGcs{
\Question{
Posloupnost je dána vzorcem pro \(n\)-tý člen ve tvaru \(a_n = \frac{n^2-4}{n+4}\).  Který člen této posloupnosti je roven \(5\)?}
{osmý}{třetí}{devátý}{pátý}{}{}}
\LANGsk{
\Question{
Postupnosť je daná vzorcom pre \(n\)-tý člen, v tvare \(a_n = \frac{n^2-4}{n+4}\).  Ktorý člen tejto postupnosti je rovný \(5\)?}
{ôsmy}{tretí}{deviaty}{piaty}{}{}}
\LANGes{
\Question{
El término general de una sucesión viene dado por \(a_n = \frac{n^2-4}{n+4}\). ¿Qué término de la sucesión equivale a \(5\)?}
{el octavo término}{el tercer término}{el noveno término}{el quinto término}{}{}}
\End

\Data\ID{2010000706}
\Updated{2022-08-25 10:08:02}
\Level{B}
\SubArea{Properties of Sequences}
\IMG{2010001401_6-figure0.pdf}
\LANGen{
\Question{
We are given a sequence \( \left( a_n \right)^{6}_{n=1} \) defined by the following graph. Find the recursive formula of such sequence.}
{\( a_1=-2, \ a_{n+1}=-a_n, \ n \in \{1;2;3;4;5\}\)}{\( a_1=2, \ a_{n+1}=-a_n, \ n \in \{1;2;3;4;5\}\)}{\( a_1=-2, \ a_{n+1}=-2a_n, \ n \in \{1;2;3;4;5\}\)}{\( a_1=2, \ a_{n+1}=a_n, \ n \in \{1;2;3;4;5\}\)}{}{}}
\LANGpl{
\Question{
Mamy ciąg \( \left( a_n \right)^{6}_{n=1} \) wyrażony na poniższym wykresie. Znajdź wzór rekurencyjny takiego ciągu.}
{\( a_1=-2, \ a_{n+1}=-a_n, \ n \in \{1;2;3;4;5\}\)}{\( a_1=2, \ a_{n+1}=-a_n, \ n \in \{1;2;3;4;5\}\)}{\( a_1=-2, \ a_{n+1}=-2a_n, \ n \in \{1;2;3;4;5\}\)}{\( a_1=2, \ a_{n+1}=a_n, \ n \in \{1;2;3;4;5\}\)}{}{}}
\LANGcs{
\Question{
Posloupnost \( \left( a_n \right)^{6}_{n=1} \) je definována grafem na obrázku. Určete rekurentní vyjádření této posloupnosti.}
{\( a_1=-2, \ a_{n+1}=-a_n, \ n \in \{1;2;3;4;5\}\)}{\( a_1=2, \ a_{n+1}=-a_n, \ n \in \{1;2;3;4;5\}\)}{\( a_1=-2, \ a_{n+1}=-2a_n, \ n \in \{1;2;3;4;5\}\)}{\( a_1=2, \ a_{n+1}=a_n, \ n \in \{1;2;3;4;5\}\)}{}{}}
\LANGsk{
\Question{
Postupnosť \( \left( a_n \right)^{6}_{n=1} \) je definovaná grafom na obrázku. Určte rekurentné vyjadrenie tejto postupnosti.}
{\( a_1=-2, \ a_{n+1}=-a_n, \ n \in \{1;2;3;4;5\}\)}{\( a_1=2, \ a_{n+1}=-a_n, \ n \in \{1;2;3;4;5\}\)}{\( a_1=-2, \ a_{n+1}=-2a_n, \ n \in \{1;2;3;4;5\}\)}{\( a_1=2, \ a_{n+1}=a_n, \ n \in \{1;2;3;4;5\}\)}{}{}}
\LANGes{
\Question{
La sucesión \( \left( a_n \right)^{6}_{n=1} \) viene dada por el gráfico en la imagen. Halla la fórmula recursiva de la sucesión.}
{\( a_1=-2, \ a_{n+1}=-a_n, \ n \in \{1;2;3;4;5\}\)}{\( a_1=2, \ a_{n+1}=-a_n, \ n \in \{1;2;3;4;5\}\)}{\( a_1=-2, \ a_{n+1}=-2a_n, \ n \in \{1;2;3;4;5\}\)}{\( a_1=2, \ a_{n+1}=a_n, \ n \in \{1;2;3;4;5\}\)}{}{}}
\End

\Data\ID{2010000707}
\Updated{2022-08-25 10:08:02}
\Level{B}
\SubArea{Properties of Sequences}
\IMG{2010001401_6-figure1.pdf}
\LANGen{
\Question{
We are given a sequence \( \left( a_n \right)^{6}_{n=1} \) defined by the following graph. Find the recursive formula of such sequence.}
{\( a_1=2, \ a_{n+1}=-a_n, \ n \in \{1;2;3;4;5\}\)}{\( a_1=-2, \ a_{n+1}=-a_n, \ n \in \{1;2;3;4;5\}\)}{\( a_1=-2, \ a_{n+1}=-2a_n, \ n \in \{1;2;3;4;5\}\)}{\( a_1=2, \ a_{n+1}=a_n, \ n \in \{1;2;3;4;5\}\)}{}{}}
\LANGpl{
\Question{
Mamy ciąg \( \left( a_n \right)^{6}_{n=1} \) przedstawiony na poniższym wykresie. Znajdź wzór rekurencyjny takiego ciągu.}
{\( a_1=2, \ a_{n+1}=-a_n, \ n \in \{1;2;3;4;5\}\)}{\( a_1=-2, \ a_{n+1}=-a_n, \ n \in \{1;2;3;4;5\}\)}{\( a_1=-2, \ a_{n+1}=-2a_n, \ n \in \{1;2;3;4;5\}\)}{\( a_1=2, \ a_{n+1}=a_n, \ n \in \{1;2;3;4;5\}\)}{}{}}
\LANGcs{
\Question{
Posloupnost \( \left( a_n \right)^{6}_{n=1} \) je definována grafem na obrázku. Určete rekurentní vyjádření této posloupnosti.}
{\( a_1=2, \ a_{n+1}=-a_n, \ n \in \{1;2;3;4;5\}\)}{\( a_1=-2, \ a_{n+1}=-a_n, \ n \in \{1;2;3;4;5\}\)}{\( a_1=-2, \ a_{n+1}=-2a_n, \ n \in \{1;2;3;4;5\}\)}{\( a_1=2, \ a_{n+1}=a_n, \ n \in \{1;2;3;4;5\}\)}{}{}}
\LANGsk{
\Question{
Postupnosť \( \left( a_n \right)^{6}_{n=1} \) je definovaná grafom na obrázku. Určte rekurentné vyjadrenie tejto postupnosti.}
{\( a_1=2, \ a_{n+1}=-a_n, \ n \in \{1;2;3;4;5\}\)}{\( a_1=-2, \ a_{n+1}=-a_n, \ n \in \{1;2;3;4;5\}\)}{\( a_1=-2, \ a_{n+1}=-2a_n, \ n \in \{1;2;3;4;5\}\)}{\( a_1=2, \ a_{n+1}=a_n, \ n \in \{1;2;3;4;5\}\)}{}{}}
\LANGes{
\Question{
La sucesión \( \left( a_n \right)^{6}_{n=1} \) viene dada por el gráfico en la imagen. Halla la fórmula recursiva de la sucesión.}
{\( a_1=2, \ a_{n+1}=-a_n, \ n \in \{1;2;3;4;5\}\)}{\( a_1=-2, \ a_{n+1}=-a_n, \ n \in \{1;2;3;4;5\}\)}{\( a_1=-2, \ a_{n+1}=-2a_n, \ n \in \{1;2;3;4;5\}\)}{\( a_1=2, \ a_{n+1}=a_n, \ n \in \{1;2;3;4;5\}\)}{}{}}
\End

\Data\ID{2010005801}
\Updated{2022-08-25 10:08:59}
\Level{B}
\SubArea{Properties of Sequences}
\LANGen{
\Question{
The value of the \(n\)th term of a sequence is given by the expression \(a^{4n}-13\). If the second term of the sequence is \(243\), which of the following is the value of \(a\)?}
{\(-2\) or \(2\)}{\(2\)}{\(4\)}{\(-2\)}{}{}}
\LANGpl{
\Question{
Wartość \(n\)-tego wyrazu ciągu określa wyrażenie \(a^{4n}-13\). Jeśli drugim wyrazem ciągu jest \(243\), który z poniższych jest wartością \(a\)?}
{\(-2\) lub \(2\)}{\(2\)}{\(4\)}{\(-2\)}{}{}}
\LANGcs{
\Question{
Hodnota \(n\)-tého členu posloupnosti je dána výrazem \(a^{4n}-13\). Pokud je druhý člen posloupnosti \(243\), které z daných čísel může být hodnotou \(a\)?}
{\(-2\) nebo \(2\)}{\(2\)}{\(4\)}{\(-2\)}{}{}}
\LANGsk{
\Question{
Hodnota \(n\)-tého člena postupnosti je daná výrazom \(a^{4n}-13\). Ak druhý člen postupnosti je \(243\), ktoré z daných výsledkov môže byť hodnota \(a\)?}
{\(-2\) alebo \(2\)}{\(2\)}{\(4\)}{\(-2\)}{}{}}
\LANGes{
\Question{
El valor del término \(n\) de una sucesión viene dado por la expresión \(a^{4n}-13\). Si el segundo término de la sucesión es \(243\), ¿cuál de los números equivale a \(a\)?}
{\(-2\) o \(2\)}{\(2\)}{\(4\)}{\(-2\)}{}{}}
\End

\Data\ID{2010005802}
\Updated{2022-08-25 10:08:59}
\Level{B}
\SubArea{Properties of Sequences}
\LANGen{
\Question{
The value of the \(n\)th term of a sequence is given by the expression \(b^{2n}-28\). If the third term of the sequence is \(701\), which of the following is the value of \(b\)?}
{\(-3\) or \(3\)}{\(3\)}{\(2\)}{\(-3\)}{}{}}
\LANGpl{
\Question{
Wartość \(n\)-tego wyrazu ciągu określa wyrażenie \(b^{2n}-28\). Jeśli trzecim wyrazem ciągu jest \(701\), który z poniższych jest wartością \(b\)?}
{\(-3\) lub \(3\)}{\(3\)}{\(2\)}{\(-3\)}{}{}}
\LANGcs{
\Question{
Hodnota \(n\)-tého členu posloupnosti je dána výrazem \(b^{2n}-28\). Pokud je třetí člen posloupnosti \(701\), které z daných čísel může být hodnotou \(b\)?}
{\(-3\) nebo \(3\)}{\(3\)}{\(2\)}{\(-3\)}{}{}}
\LANGsk{
\Question{
Hodnota \(n\)-tého člena postupnosti je daná výrazom \(b^{2n}-28\). Ak tretí člen postupnosti je \(701\), ktoré z daných výsledkov môže byť hodnota \(b\)?}
{\(-3\) alebo \(3\)}{\(3\)}{\(2\)}{\(-3\)}{}{}}
\LANGes{
\Question{
El valor del término \(n\) de una sucesión viene dado por la expresión \(b^{2n}-28\). Si el tercer término de la sucesión es \(701\), ¿cuál de los números equivale a \(b\)?}
{\(-3\) o \(3\)}{\(3\)}{\(2\)}{\(-3\)}{}{}}
\End

\Data\ID{2010005803}
\Updated{2022-08-25 10:08:59}
\Level{B}
\SubArea{Properties of Sequences}
\IMG{2010005801-figure0.pdf}
\LANGen{
\Question{
We are given a sequence \( (2+(-1)^n)^{\infty}_{n=1}\). A part of the sequence is displayed in the graph. Complete the sentence: This sequence is ...}
{neither increasing nor decreasing.}{increasing.}{decreasing.}{not bounded.}{}{}}
\LANGpl{
\Question{
Dany jest ciąg \( (2+(-1)^n)^{\infty}_{n=1}\). Część ciągu jest przedstawiana na wykresie. Uzupełnij zdanie: Ten ciąg}
{nie jest ani rosnący, ani malejący.}{jest rosnący}{jest malejący.}{nie jest ograniczony.}{}{}}
\LANGcs{
\Question{
Na obrázku je část grafu posloupnosti \( (2+(-1)^n)^{\infty}_{n=1}\). Doplňte výrok: Posloupnost ...}
{není ani roustoucí, ani klesající.}{je rostoucí.}{je klesající.}{není omezená.}{}{}}
\LANGsk{
\Question{
Daná je postupnosť \( (2+(-1)^n)^{\infty}_{n=1}\), ktorej časť grafu je na obrázku. Dokončite vetu: Táto postupnosť ...}
{nie je rastúca, ani klesajúca.}{je rastúca.}{je klesajúca.}{nie je ohraničená.}{}{}}
\LANGes{
\Question{
Dada la sucesión \( (2+(-1)^n)^{\infty}_{n=1}\). Una parte de la sucesión la representa el siguiente gráfico. Se trata de una sucesión...}
{ni creciente ni decreciente.}{creciente.}{decreciente.}{no acotada.}{}{}}
\End

\Data\ID{2010005804}
\Updated{2022-08-25 10:08:59}
\Level{B}
\SubArea{Properties of Sequences}
\IMG{2010005801-figure1.pdf}
\LANGen{
\Question{
We are given a sequence \( \left(\sin \left(n\frac{\pi}{2}\right)\right)^{\infty}_{n=1}\). A part of the sequence is displayed in the graph. Complete the sentence: This sequence is ...}
{neither increasing nor decreasing.}{increasing.}{decreasing.}{not bounded.}{}{}}
\LANGpl{
\Question{
Mamy dany ciąg \( \left(\sin \left(n\frac{\pi}{2}\right)\right)^{\infty}_{n=1}\). Część ciągu jest przedstawiana na wykresie. Uzupełnij zdanie: Ten ciąg ...}
{nie jest rosnący ani malejący.}{jest rosnący.}{jest malejący.}{nie jest ograniczony.}{}{}}
\LANGcs{
\Question{
Na obrázku je část grafu posloupnosti \( \left(\sin \left(n\frac{\pi}{2}\right)\right)^{\infty}_{n=1}\). Doplňte výrok: Posloupnost  ...}
{není ani roustoucí, ani klesající.}{je rostoucí.}{je klesající.}{není omezená.}{}{}}
\LANGsk{
\Question{
Daná je postupnosť \( \left(\sin \left(n\frac{\pi}{2}\right)\right)^{\infty}_{n=1}\), ktorej časť grafu je na obrázku. Dokončite vetu: Táto postupnosť ...}
{nie je rastúca, ani klesajúca.}{je rastúca.}{je klesajúca.}{nie je ohraničená.}{}{}}
\LANGes{
\Question{
Dada la sucesión \( \left(\sin \left(n\frac{\pi}{2}\right)\right)^{\infty}_{n=1}\). Una parte de la sucesión la representa el siguiente gráfico. Se trata de una sucesión ...}
{ni creciente ni decreciente.}{creciente.}{decreciente.}{no acotada.}{}{}}
\End

\Data\ID{2010010301}
\Updated{2022-11-02 09:11:35}
\Level{B}
\SubArea{Properties of Sequences}
\LANGen{
\Question{
Which of the following sequences given by the formula for the \(n\)th term is not decreasing?}
{\( (\log n)^{\infty}_{n=1}\)}{\( (4-n^2)^{\infty}_{n=1}\)}{\( \left( \frac{2n+1}{n}\right)^{\infty}_{n=1}\)}{\( \left( \frac{1}{2^n}\right)^{\infty}_{n=1}\)}{}{}}
\LANGpl{
\Question{
Który z poniższych ciągów podanych za pomocą wzoru na \(n\)-ty wyraz nie maleje?}
{\( (\log n)^{\infty}_{n=1}\)}{\( (4-n^2)^{\infty}_{n=1}\)}{\( \left( \frac{2n+1}{n}\right)^{\infty}_{n=1}\)}{\( \left( \frac{1}{2^n}\right)^{\infty}_{n=1}\)}{}{}}
\LANGcs{
\Question{
Která z posloupností zadaná vzorcem pro \(n\)-tý člen není klesající?}
{\( (\log n)^{\infty}_{n=1}\)}{\( (4-n^2)^{\infty}_{n=1}\)}{\( \left( \frac{2n+1}{n}\right)^{\infty}_{n=1}\)}{\( \left( \frac{1}{2^n}\right)^{\infty}_{n=1}\)}{}{}}
\LANGsk{
\Question{
Ktorá postupnosť daná vzťahom pre \(n\)-tý člen nie je klesajúca?}
{\( (\log n)^{\infty}_{n=1}\)}{\( (4-n^2)^{\infty}_{n=1}\)}{\( \left( \frac{2n+1}{n}\right)^{\infty}_{n=1}\)}{\( \left( \frac{1}{2^n}\right)^{\infty}_{n=1}\)}{}{}}
\LANGes{
\Question{
¿Cuál de las siguiente sucesiones dadas por la fórmula del término general no es decreciente?}
{\( (\log n)^{\infty}_{n=1}\)}{\( (4-n^2)^{\infty}_{n=1}\)}{\( \left( \frac{2n+1}{n}\right)^{\infty}_{n=1}\)}{\( \left( \frac{1}{2^n}\right)^{\infty}_{n=1}\)}{}{}}
\End

\Data\ID{2010010302}
\Updated{2022-11-02 09:11:35}
\Level{B}
\SubArea{Properties of Sequences}
\LANGen{
\Question{
Which of the following sequences given by the formula for the \(n\)th term is not increasing?}
{\( (n^{-4})^{\infty}_{n=1}\)}{\( (\sqrt{n})^{\infty}_{n=1}\)}{\( \left( -\frac{n+1}{n}\right)^{\infty}_{n=1}\)}{\( \left( {2^n}\right)^{\infty}_{n=1}\)}{}{}}
\LANGpl{
\Question{
Który z następujących ciągów podanych za pomocą wzoru na \(n\)-ty wyraz nie rośnie?}
{\( (n^{-4})^{\infty}_{n=1}\)}{\( (\sqrt{n})^{\infty}_{n=1}\)}{\( \left( -\frac{n+1}{n}\right)^{\infty}_{n=1}\)}{\( \left( {2^n}\right)^{\infty}_{n=1}\)}{}{}}
\LANGcs{
\Question{
Která z posloupností zadaná vzorcem pro \(n\)-tý člen není rostoucí?}
{\( (n^{-4})^{\infty}_{n=1}\)}{\( (\sqrt{n})^{\infty}_{n=1}\)}{\( \left( -\frac{n+1}{n}\right)^{\infty}_{n=1}\)}{\( \left( {2^n}\right)^{\infty}_{n=1}\)}{}{}}
\LANGsk{
\Question{
Ktorá postupnosť daná vzťahom pre \(n\)-tý člen nie je rastúca?}
{\( (n^{-4})^{\infty}_{n=1}\)}{\( (\sqrt{n})^{\infty}_{n=1}\)}{\( \left( -\frac{n+1}{n}\right)^{\infty}_{n=1}\)}{\( \left( {2^n}\right)^{\infty}_{n=1}\)}{}{}}
\LANGes{
\Question{
¿Cuál de las siguiente sucesiones dadas por la fórmula del término general no es creciente?}
{\( (n^{-4})^{\infty}_{n=1}\)}{\( (\sqrt{n})^{\infty}_{n=1}\)}{\( \left( -\frac{n+1}{n}\right)^{\infty}_{n=1}\)}{\( \left( {2^n}\right)^{\infty}_{n=1}\)}{}{}}
\End

\Data\ID{2010010303}
\Updated{2022-11-02 09:11:35}
\Level{B}
\SubArea{Properties of Sequences}
\LANGen{
\Question{
We are given a sequence \( \left( \frac{2n+1}{n+3}\right)^{\infty}_{n=1}\). What are the properties of this sequence?}
{increasing and bounded}{neither increasing nor decreasing}{decreasing and bounded above}{increasing and not bounded above}{decreasing and not bounded above}{}}
\LANGpl{
\Question{
Otrzymujemy ciąg \( \left( \frac{2n+1}{n+3}\right)^{\infty}_{n=1}\). Jakie są właściwości tego ciągu?}
{rosnący i ograniczony}{ani rosnący ani malejący}{malejący i ograniczony z góry}{rosnący i nie ograniczony z dołu}{malejący i nie ograniczony z góry}{}}
\LANGcs{
\Question{
Je dána posloupnost \( \left( \frac{2n+1}{n+3}\right)^{\infty}_{n=1}\). Určete vlastnosti dané posloupnosti.}
{je roustoucí a omezená}{není ani rostoucí, ani klesající}{je klesající a omezená shora}{je rostoucí a není omezená shora}{je klesající a není omezená shora}{}}
\LANGsk{
\Question{
Daná je postupnosť \( \left( \frac{2n+1}{n+3}\right)^{\infty}_{n=1}\). Označte vlastnosti tejto postupnosti.}
{rastúca a ohraničená}{nie je rastúca ani klesajúca}{klesajúca a zhora ohraničená}{rastúca a nie je zhora ohraničená}{klesajúca a nie je zhora ohraničená}{}}
\LANGes{
\Question{
Dada la sucesión \( \left( \frac{2n+1}{n+3}\right)^{\infty}_{n=1}\). ¿Cuáles son las propiedades de esta sucesión?}
{creciente y acotada}{ni creciente ni decreciente}{decreciente y acotada superiormente}{creciente y no acotada superiormente}{decreciente y no acotada superiormente}{}}
\End

\Data\ID{2010010304}
\Updated{2022-11-02 09:11:50}
\Level{B}
\SubArea{Properties of Sequences}
\LANGen{
\Question{
We are given a sequence \( \left( \frac{n+5}{n+1}\right)^{\infty}_{n=1}\). What are the properties of this sequence?}
{decreasing and bounded above}{increasing and not bounded below}{neither increasing nor decreasing}{increasing and bounded below}{decreasing and not bounded above}{}}
\LANGpl{
\Question{
Otrzymujemy ciąg \( \left( \frac{n+5}{n+1}\right)^{\infty}_{n=1}\). Jakie są właściwości tego ciągu?}
{malejący i ograniczony z góry}{rosnący i nie ograniczony z dołu}{ani rosnący ani malejący}{rosnący i ograniczony z dołu}{malejący i nie ograniczony z góry}{}}
\LANGcs{
\Question{
Je dána posloupnost \( \left( \frac{n+5}{n+1}\right)^{\infty}_{n=1}\). Určete její vlastnosti.}
{je klesající a omezená shora}{je rostoucí a není omezená zdola}{není ani rostoucí, ani klesající}{je rostoucí a omezená zdola}{je klesající a není omezená zdola}{}}
\LANGsk{
\Question{
Daná je postupnosť \( \left( \frac{n+5}{n+1}\right)^{\infty}_{n=1}\). Označte vlastnosti tejto postupnosti.}
{klesajúca a zhora ohraničená}{rastúca a nie je zdola ohraničená}{nie je rastúca ani klesajúca}{rastúca a zdola ohraničená}{klesajúca a nie je zhora ohraničená}{}}
\LANGes{
\Question{
Dada la sucesión \( \left( \frac{n+5}{n+1}\right)^{\infty}_{n=1}\). ¿Cuáles son las propiedades de esta sucesión?}
{decreciente y acotada superiormente}{creciente y no acotada inferiormente}{ni creciente ni decreciente}{creciente y acotada inferiormente}{decreciente y no acotada superiormente}{}}
\End

\Data\ID{2010010305}
\Updated{2023-03-27 08:03:51}
\Level{B}
\SubArea{Properties of Sequences}
\LANGen{
\Question{
Which of following statements about the sequence \( \left( \frac{n+3}{2n}\right)^{\infty}_{n=1} \) is true? 
\[\]
(Help: A sequence is bounded below if all its terms are greater than or equal to a real number \(L\), which is called the lower bound of the sequence. Similarly, a sequence is bounded above if all its terms are less than or equal to a real number \(U\), which is called the upper bound of the sequence.)}
{one of lower bounds is \(\frac12\), one of upper bounds is \(2\)}{one of lower bounds is \(\frac12\), upper bound does not exist}{lower bound does not exist, one of upper bounds is \(2\)}{there exists neither lower bound nor upper bound}{}{}}
\LANGpl{
\Question{
Które z poniższych stwierdzeń o ciągu \( \left( \frac{n+3}{2n}\right)^{\infty}_{n=1} \) jest prawdziwe?
\[\]
(Wskazówka: ciąg jest ograniczony z dołu, jeśli wszystkie jego wyrazy są większe lub równe liczbie rzeczywistej \(L\), która jest nazywana dolną granicą ciągu. Podobnie ciąg jest ograniczony z góry, jeśli wszystkie jego wyrazy są mniejsze lub równe liczbie rzeczywistej \(U\), która jest nazywana górną granicą ciągu.}
{jedną z dolnych granic jest  \(\frac12\),   jedną z górnych granic jest \(2\)}{jedną z dolnych granic jest \(\frac12\), górna granica nie istnieje}{dolna granica nie istnieje,  jedną z górnych granic jest \(2\)}{nie istnieje ani górna ani dolna granica}{}{}}
\LANGcs{
\Question{
Které z následujících tvrzení o poslounosti \( \left( \frac{n+3}{2n}\right)^{\infty}_{n=1} \) je pravdivé? 
\[\]
(Nápověda: Posloupnost je omezená zdola, pokud jsou všechny její členy větší nebo rovny reálnému číslu \(L\), které nazýváme dolní mez posloupnosti. Podobně, posloupnost je omezená shora, pokud jsou všechny její členy menší nebo rovny reálnému číslu \(U\), které nazýváme horní mez posloupnosti.)}
{jedna z dolních mezí je \(\frac12\), jedna z horních mezí je \(2\)}{jedna z dolních mezí je \(\frac12\), horní mez neexistuje}{dolní mez neexistuje, jedna z horních mezí je \(2\)}{neexistuje ani dolní ani horní mez}{}{}}
\LANGsk{
\Question{
Ktoré z nasledujúcich tvrdení o postupnosti \( \left( \frac{n+3}{2n}\right)^{\infty}_{n=1} \) je pravdivé. 
\[\]
(Pomôcka: Postupnosť je zdola ohraničená, ak všetky jej členy sú väčšie alebo rovné reálnemu číslu \(d\), ktorému hovoríme dolné ohraničenie postupnosti. Postupnosť je zhora ohraničená, ak všetky jej členy sú menšie alebo rovné reálnemu číslu \(h\), ktorému hovoríme horné ohraničenie postupnosti.)}
{Jedna z dolných hraníc je \(\frac12\) a jedna z horných hraníc je \(2\).}{Jedna z dolných hraníc je \(\frac12\) a horná hranica neexistuje.}{Dolná hranica neexistuje a jedna z horných hraníc je \(2\).}{Neexistuje ani dolná ani horná hranica.}{}{}}
\LANGes{
\Question{
¿Cuál de las siguientes proposiciones lógicas sobre esta sucesión \( \left( \frac{n+3}{2n}\right)^{\infty}_{n=1} \) es verdadera? 
\[\]

(Sugerencia: una sucesión está acotada por abajo si todos sus términos son mayores o iguales a un número real \(L\), lo que se denomina límite inferior de la sucesión. Así, una sucesión está acotada por arriba si todos sus números son menores o iguales a un número real \(U\), lo que se denomina límite superior de la sucesión).}
{uno de los límites inferiores es \(\frac12\), uno de los límites superiores es \(2\)}{uno de los límites inferiores es \(\frac12\), el límite superior no existe}{el límite inferior no existe, uno de los límites superiores es \(2\)}{no existe límite inferior ni límite superior}{}{}}
\End

\Data\ID{2010010306}
\Updated{2023-03-27 08:03:19}
\Level{B}
\SubArea{Properties of Sequences}
\LANGen{
\Question{
Which of following statements about the sequence \( \left( \frac{n-2}{n+1}\right)^{\infty}_{n=1} \) is true? 
\[\]
(Help: A sequence is bounded below if all its terms are greater than or equal to a real number \(L\), which is called the lower bound of the sequence. Similarly, a sequence is bounded above if all its terms are less than or equal to a real number \(U\), which is called the upper bound of the sequence.)}
{one of lower bounds is \(-\frac12\), one of upper bounds is \(1\)}{one of lower bounds is \(-\frac12\), upper bound does not exist}{lower bound does not exist, one of upper bounds is \(1\)}{there exists neither lower bound nor upper bound}{}{}}
\LANGpl{
\Question{
Które z poniższych stwierdzeń o ciągu \( \left( \frac{n-2}{n+1}\right)^{\infty}_{n=1} \) jest prawdziwe?
\[\]
(Wskazówka: ciąg jest ograniczony z dołu, jeśli wszystkie jej wyrazy są większe lub równe liczbie rzeczywistej \(L\), która jest nazywana dolną granicą ciągu. Podobnie ciąg jest ograniczony z góry, jeśli wszystkie jego wyrazy są mniejsze lub równe liczbie rzeczywistej \(U\), która jest nazywana górną granicą ciągu.)}
{jedna z dolnych granic jest  \(-\frac12\), jedną z górnych granic jest \(1\)}{jedna z dolnych granic jest  \(-\frac12\), górna granica nie istnieje}{dolna granica nie istnieje, jedna z górnych granic jest równa \(1\)}{nie istnieje ani dolna ani górna granica}{}{}}
\LANGcs{
\Question{
Které z následujících tvrzení o poslounosti \( \left( \frac{n-2}{n+1}\right)^{\infty}_{n=1} \) je pravdivé? 
\[\]
(Nápověda: Posloupnost je omezená zdola, pokud jsou všechny její členy větší nebo rovny reálnému číslu \(L\), které nazýváme dolní mez posloupnosti. Podobně, posloupnost je omezená shora, pokud jsou všechny její členy menší nebo rovny reálnému číslu \(U\), které nazýváme horní mez posloupnosti.)}
{jedna z dolních mezí je \(-\frac12\), jedna z horních mezí je \(1\)}{jedna z dolních mezí je \(-\frac12\), horní mez neexistuje}{dolní mez neexistuje, jedna z horních mezí je \(1\)}{neexistuje ani horní ani dolní mez}{}{}}
\LANGsk{
\Question{
Ktoré z nasledujúcich tvrdení o postupnosti \( \left( \frac{n-2}{n+1}\right)^{\infty}_{n=1} \) je pravdivé.
\[\]
(Pomôcka: Postupnosť je zdola ohraničená, ak všetky jej členy sú väčšie alebo rovné reálnemu číslu \(d\), ktorému hovoríme dolné ohraničenie postupnosti. Postupnosť je zhora ohraničená, ak všetky jej členy sú menšie alebo rovné reálnemu číslu \(h\), ktorému hovoríme horné ohraničenie postupnosti.)}
{Jedna z dolných hraníc je \(-\frac12\) a jedna z horných hraníc je \(1\).}{Jedna z dolných hraníc je \(-\frac12\) a horná hranica neexistuje.}{Dolná hranica neexistuje a jedna z horných hraníc je \(1\).}{Neexistuje ani dolná ani horná hranica.}{}{}}
\LANGes{
\Question{
¿Cuál de las siguientes proposiciones lógicas sobre esta sucesión \( \left( \frac{n-2}{n+1}\right)^{\infty}_{n=1} \) es verdadera?
\[\]

(Sugerencia: una sucesión está acotada por abajo si todos sus términos son mayores o iguales a un número real \(L\), lo que se denomina límite inferior de la sucesión. Así, una sucesión está acotada por arriba si todos sus números son menores o iguales a un número real \(U\), lo que se denomina límite superior de la sucesión).}
{uno de los límites inferiores es \(-\frac12\), uno de los límites superiores es \(1\)}{uno de los límites inferiores es \(-\frac12\), el límite superior no existe}{el límite inferior no existe, uno de los límites superiores es \(1\)}{no existe límite inferior ni límite superior}{}{}}
\End

\Data\ID{9000063801}
\Updated{2020-07-15 03:07:37}
\Level{B}
\SubArea{Properties of Sequences}
\LANGen{
\Question{
We are given the sequence \(\left (an + b\right )_{n=1}^{\infty }\).
This sequence satisfies \(a_{2} = 2\) 
and \(a_{4} = 8\). Use this
information to find \(a\).}
{\(a = 3\)}{\(a = 1\)}{\(a = 2\)}{\(a = 4\)}{}{}}
\LANGpl{
\Question{
Dany jest ciąg  \(\left (an + b\right )_{n=1}^{\infty }\).
Ciąg spełnia  \(a_{2} = 2\) 
i \(a_{4} = 8\). Wykorzystaj te informacje, aby wyznaczyć \(a\).}
{\(a = 3\)}{\(a = 1\)}{\(a = 2\)}{\(a = 4\)}{}{}}
\LANGcs{
\Question{
Je dána posloupnost \(\left (an + b\right )_{n=1}^{\infty }\),
ve které platí, že \(a_{2} = 2\) 
a \(a_{4} = 8\).
Najděte $a$.}
{\(a = 3\)}{\(a = 1\)}{\(a = 2\)}{\(a = 4\)}{}{}}
\LANGsk{
\Question{
Je daná postupnosť \(\left (an + b\right )_{n=1}^{\infty }\),
v ktorej platí, že \(a_{2} = 2\) 
a \(a_{4} = 8\).
Nájdite $a$.}
{\(a = 3\)}{\(a = 1\)}{\(a = 2\)}{\(a = 4\)}{}{}}
\LANGes{
\Question{
Dada la sucesión \(\left (an + b\right )_{n=1}^{\infty }\) en la que vale \(a_{2} = 2\) 
y \(a_{4} = 8\). Halla \(a\).}
{\(a = 3\)}{\(a = 1\)}{\(a = 2\)}{\(a = 4\)}{}{}}
\End

\Data\ID{9000063802}
\Updated{2022-04-23 05:04:02}
\Level{B}
\SubArea{Properties of Sequences}
\LANGen{
\Question{
We are given the sequence \(\left (an + b\right )_{n=1}^{\infty }\).
This sequence satisfies \(a_{4} - a_{1} = 6\). Use
this information to find \(a\).}
{\(a = 2\)}{\(a = -2\)}{\(a = -1\)}{\(a = 1\)}{}{}}
\LANGpl{
\Question{
Dany jest ciąg  \(\left (an + b\right )_{n=1}^{\infty }\), który spełnia  \(a_{4} - a_{1} = 6\). Wykorzystaj te informacje, by wyznaczyć 
\(a\).}
{\(a = 2\)}{\(a = -2\)}{\(a = -1\)}{\(a = 1\)}{}{}}
\LANGcs{
\Question{
Je dána posloupnost \(\left (an + b\right )_{n=1}^{\infty }\),
ve které platí, že \(a_{4} - a_{1} = 6\).
Najděte $a$.}
{\(a = 2\)}{\(a = -2\)}{\(a = -1\)}{\(a = 1\)}{}{}}
\LANGsk{
\Question{
Je daná postupnosť \(\left (an + b\right )_{n=1}^{\infty }\),
v ktorej platí, že \(a_{4} - a_{1} = 6\). Nájdite $a$.}
{\(a = 2\)}{\(a = -2\)}{\(a = -1\)}{\(a = 1\)}{}{}}
\LANGes{
\Question{
Dada la sucesión \(\left (an + b\right )_{n=1}^{\infty }\) en la que vale \(a_{4} - a_{1} = 6\). Halla \(a\).}
{\(a = 2\)}{\(a = -2\)}{\(a = -1\)}{\(a = 1\)}{}{}}
\End

\Data\ID{9000063806}
\Updated{2023-03-27 08:03:57}
\Level{B}
\SubArea{Properties of Sequences}
\LANGen{
\Question{
Consider the sequence \(a_{n+1} = a_{n} - 2a_{n-1}\) 
with \(a_{3} = 0\) 
and \(a_{4} = -16\).
Find \(a_{2} - a_{1}\).}
{\(4\)}{\(16\)}{\(- 4\)}{\(8\)}{}{}}
\LANGpl{
\Question{
Rozważ ciąg  \(a_{n+1} = a_{n} - 2a_{n-1}\) 
z \(a_{3} = 0\) 
i \(a_{4} = -16\).
Oblicz \(a_{2} - a_{1}\).}
{\(4\)}{\(16\)}{\(- 4\)}{\(8\)}{}{}}
\LANGcs{
\Question{
Uvažujme rekurentně zadanou posloupnost
\(a_{n+1} = a_{n} - 2a_{n-1}\), kde
\(a_{3} = 0\) a
\(a_{4} = -16\).
Potom platí:}
{\(a_{2} - a_{1} = 4\)}{\(a_{2} - a_{1} = 16\)}{\(a_{2} - a_{1} = -4\)}{\(a_{2} - a_{1} = 8\)}{}{}}
\LANGsk{
\Question{
Je daná rekurentne zadaná postupnosť
\(a_{n+1} = a_{n} - 2a_{n-1}\), kde
\(a_{3} = 0\) a
\(a_{4} = -16\).
Potom platí:}
{\(a_{2} - a_{1} = 4\)}{\(a_{2} - a_{1} = 16\)}{\(a_{2} - a_{1} = -4\)}{\(a_{2} - a_{1} = 8\)}{}{}}
\LANGes{
\Question{
Dada la sucesión \(a_{n+1} = a_{n} - 2a_{n-1}\) 
con \(a_{3} = 0\) 
y \(a_{4} = -16\).
Halla \(a_{2} - a_{1}\).}
{\(4\)}{\(16\)}{\(- 4\)}{\(8\)}{}{}}
\End

\Data\ID{9000063808}
\Updated{2020-07-15 03:07:37}
\Level{B}
\SubArea{Properties of Sequences}
\LANGen{
\Question{
Given the sequence \(\left (2n + 3\right )_{n=1}^{\infty }\),
find the recurrence relation for this sequence.}
{\(a_{n+1} = a_{n} + 2,\ a_{1} = 5\)}{\(a_{n+1} = a_{n} + 3,\ a_{1} = 5\)}{\(a_{n+1} = a_{n} + 4,\ a_{1} = 5\)}{\(a_{n+1} = a_{n} + 5,\ a_{1} = 5\)}{}{}}
\LANGpl{
\Question{
Dany jest ciąg  \(\left (2n + 3\right )_{n=1}^{\infty }\). Określ relację rekurencyjną tego ciągu.}
{\(a_{n+1} = a_{n} + 2,\ a_{1} = 5\)}{\(a_{n+1} = a_{n} + 3,\ a_{1} = 5\)}{\(a_{n+1} = a_{n} + 4,\ a_{1} = 5\)}{\(a_{n+1} = a_{n} + 5,\ a_{1} = 5\)}{}{}}
\LANGcs{
\Question{
Je dána posloupnost \(\left (2n + 3\right )_{n=1}^{\infty }\).
Rekurentní vyjádření této posloupnosti je:}
{\(a_{n+1} = a_{n} + 2,\ a_{1} = 5\)}{\(a_{n+1} = a_{n} + 3,\ a_{1} = 5\)}{\(a_{n+1} = a_{n} + 4,\ a_{1} = 5\)}{\(a_{n+1} = a_{n} + 5,\ a_{1} = 5\)}{}{}}
\LANGsk{
\Question{
Je daná postupnosť \(\left (2n + 3\right )_{n=1}^{\infty }\).
Rekurentné vyjadrenie tejto postupnosti je:}
{\(a_{n+1} = a_{n} + 2,\ a_{1} = 5\)}{\(a_{n+1} = a_{n} + 3,\ a_{1} = 5\)}{\(a_{n+1} = a_{n} + 4,\ a_{1} = 5\)}{\(a_{n+1} = a_{n} + 5,\ a_{1} = 5\)}{}{}}
\LANGes{
\Question{
Dada la sucesión \(\left (2n + 3\right )_{n=1}^{\infty }\), Halla la fórmula recursiva de la sucesión.}
{\(a_{n+1} = a_{n} + 2,\ a_{1} = 5\)}{\(a_{n+1} = a_{n} + 3,\ a_{1} = 5\)}{\(a_{n+1} = a_{n} + 4,\ a_{1} = 5\)}{\(a_{n+1} = a_{n} + 5,\ a_{1} = 5\)}{}{}}
\End

\Data\ID{9000063809}
\Updated{2020-07-15 03:07:37}
\Level{B}
\SubArea{Properties of Sequences}
\LANGen{
\Question{
Given the sequence \(\left ( \frac{1}
{n(n+1)}\right )_{n=1}^{\infty }\),
find the recurrence relation for this sequence.}
{\(a_{n+1} = \frac{n} 
{n+2}a_{n},\ a_{1} = \frac{1} 
{2}\)}{\(a_{n+1} = \frac{n} 
{n+1}a_{n},\ a_{1} = \frac{1} 
{2}\)}{\(a_{n+1} = \frac{n+1} 
 {n} a_{n},\ a_{1} = \frac{1} 
{2}\)}{\(a_{n+1} = \frac{n+1} 
{n+2}a_{n},\ a_{1} = \frac{1} 
{2}\)}{}{}}
\LANGpl{
\Question{
Dany jest ciąg \(\left ( \frac{1}
{n(n+1)}\right )_{n=1}^{\infty }\).
Wyznacz relację rekurencyjną tego ciągu.}
{\(a_{n+1} = \frac{n} 
{n+2}a_{n},\ a_{1} = \frac{1} 
{2}\)}{\(a_{n+1} = \frac{n} 
{n+1}a_{n},\ a_{1} = \frac{1} 
{2}\)}{\(a_{n+1} = \frac{n+1} 
 {n} a_{n},\ a_{1} = \frac{1} 
{2}\)}{\(a_{n+1} = \frac{n+1} 
{n+2}a_{n},\ a_{1} = \frac{1} 
{2}\)}{}{}}
\LANGcs{
\Question{
Je dána posloupnost \(\left ( \frac{1}
{n(n+1)}\right )_{n=1}^{\infty }\).
Rekurentní vyjádření této posloupnosti je:}
{\(a_{n+1} = \frac{n} 
{n+2}a_{n},\ a_{1} = \frac{1} 
{2}\)}{\(a_{n+1} = \frac{n} 
{n+1}a_{n},\ a_{1} = \frac{1} 
{2}\)}{\(a_{n+1} = \frac{n+1} 
 {n} a_{n},\ a_{1} = \frac{1} 
{2}\)}{\(a_{n+1} = \frac{n+1} 
{n+2}a_{n},\ a_{1} = \frac{1} 
{2}\)}{}{}}
\LANGsk{
\Question{
Je daná postupnosť \(\left ( \frac{1}
{n(n+1)}\right )_{n=1}^{\infty }\).
Rekurentné vyjadrenie tejto postupnosti je:}
{\(a_{n+1} = \frac{n} 
{n+2}a_{n},\ a_{1} = \frac{1} 
{2}\)}{\(a_{n+1} = \frac{n} 
{n+1}a_{n},\ a_{1} = \frac{1} 
{2}\)}{\(a_{n+1} = \frac{n+1} 
 {n} a_{n},\ a_{1} = \frac{1} 
{2}\)}{\(a_{n+1} = \frac{n+1} 
{n+2}a_{n},\ a_{1} = \frac{1} 
{2}\)}{}{}}
\LANGes{
\Question{
Dada la sucesión \(\left ( \frac{1}
{n(n+1)}\right )_{n=1}^{\infty }\),
Halla la fórmula recursiva de la sucesión.}
{\(a_{n+1} = \frac{n} 
{n+2}a_{n},\ a_{1} = \frac{1} 
{2}\)}{\(a_{n+1} = \frac{n} 
{n+1}a_{n},\ a_{1} = \frac{1} 
{2}\)}{\(a_{n+1} = \frac{n+1} 
 {n} a_{n},\ a_{1} = \frac{1} 
{2}\)}{\(a_{n+1} = \frac{n+1} 
{n+2}a_{n},\ a_{1} = \frac{1} 
{2}\)}{}{}}
\End
